\chapter{Symmetric Spaces and Holonomy}
\label{chap:pet-symmetric-spaces-and-holonomy}

Faithful transcription of Petersen, \emph{Riemannian Geometry} (GTM 171, 3rd ed.),
Chapter 10. The chapter opens (unnumbered, no formal statements) with a
motivational overview of symmetric spaces and holonomy — explicit examples such
as complex projective space, the de Rham decomposition theorem, the
relationship between holonomy and symmetric spaces, the classification of
compact manifolds with nonnegative curvature operator, and historical remarks
on Yau's resolution of the Calabi conjecture and D. Joyce's construction of
manifolds with exceptional holonomy. The formal development begins in §10.1,
which defines a symmetric space $(M,g)$ as a Riemannian manifold each of whose
points $p$ is an isolated fixed point of an isometric involution $A_p$ (with
$DA_p|_p=-I$), tabulates the classical families of compact and noncompact
symmetric spaces and introduces the notion of rank (§10.1.1); shows that
having parallel curvature tensor almost characterizes symmetric spaces, via
Cartan's two extension theorems (§10.1.2); and gives an algebraic
(Lie-theoretic) description of the isotropy representation of a symmetric
space, including the short exact sequence
$0\to\mathrm{iso}_p\to\mathrm{iso}\to T_pM\to0$ of Killing-field Lie algebras
(Petersen's Proposition 8.1.4, from outside this chapter) and a curvature
formula for the Lie bracket of $\mathrm{iso}$ (§10.1.3). The chapter goes on to
work this algebraic description out for the Grassmannians, complex projective
space, $\mathrm{SL}(n)/\mathrm{SO}(n)$, and Lie groups with biinvariant metrics
(§10.2), and finally develops the theory of the holonomy group, the de Rham
splitting theorem, and the Berger–Simons classification of holonomy groups
together with its consequences for rank and curvature rigidity (§10.3).

\section{Symmetric Spaces}

\subsection{The Homogeneous Description}

\begin{definition}[symmetric space]
  \label{def:pet-ch10-symmetric-space}
  \lean{PetersenLib.IsSymmetricSpace}
  A Riemannian manifold $(M,g)$ is a \emph{symmetric space} if for each
  $p\in M$ the isotropy group $\mathrm{Iso}_p$ contains an isometry $A_p$ (the
  \emph{point symmetry} at $p$) such that $DA_p:T_pM\to T_pM$ is the antipodal
  map $-I$. Since isometries preserve geodesics, any geodesic $c(t)$ with
  $c(0)=p$ satisfies $A_p(c(t))=c(-t)$.
  \source{petersen:page-0379}
\end{definition}

\begin{proposition}[symmetric spaces are homogeneous, hence complete]
  \label{prop:pet-ch10-symmetric-homogeneous}
  \lean{PetersenLib.symmetricSpace_isHomogeneous}
  \uses{def:pet-ch10-symmetric-space}
  A symmetric space is homogeneous, and consequently complete.
  \source{petersen:page-0379}
\end{proposition}
\begin{proof}
  \uses{prop:pet-ch10-symmetric-homogeneous}
  If two points $q,r\in M$ are joined by a geodesic, the point symmetry $A_m$
  at the midpoint $m$ of the geodesic segment between $q$ and $r$ is an
  isometry that interchanges $q$ and $r$ (parametrize the geodesic so that
  $m$ corresponds to $t=0$; then $A_m(c(t))=c(-t)$ sends the parameter value
  of $q$ to that of $r$). Since any two points of the connected manifold $M$
  can be joined by a broken sequence of geodesic segments, composing the
  corresponding midpoint symmetries carries any $p\in M$ to any $q\in M$ by an
  isometry, so $\mathrm{Iso}(M,g)$ acts transitively and $(M,g)$ is
  homogeneous. Homogeneity forces the completeness of $(M,g)$: geodesics from
  any point extend at least as far as from any other, so no geodesic can fail
  to extend for all time.
\end{proof}

\begin{remark}[a symmetry at one point suffices; biinvariant Lie groups]
  \label{rem:pet-ch10-single-point-symmetry}
  \lean{PetersenLib.symmetricSpace_of_symmetry_at_one_point,PetersenLib.biinvariantLieGroupSymmetricSpace}
  \uses{def:pet-ch10-symmetric-space}
  A homogeneous space $G/H=\mathrm{Iso}/\mathrm{Iso}_p$ is symmetric as soon as
  the point symmetry $A_p$ exists at a single point $p$: for $q=h\cdot p$ with
  $h\in\mathrm{Iso}$, $A_q:=h\circ A_p\circ h^{-1}$ is a point symmetry at $q$.
  In particular, every Lie group $G$ with a biinvariant metric is a symmetric
  space, since $g\mapsto g^{-1}$ is the desired point symmetry at the
  identity element.
  \source{petersen:page-0379}
\end{remark}

\begin{definition}[Tables 10.1--10.4 — compact groups, their noncompact duals, and the classical rank-one symmetric spaces]
  \label{def:pet-ch10-symmetric-space-tables}
  \lean{PetersenLib.symmetricSpaceTables}
  \uses{rem:pet-ch10-single-point-symmetry}
  The following families of homogeneous spaces are symmetric, always occurring
  in dual pairs of compact and noncompact spaces (there are many further
  families and several exceptional examples).
  \begin{center}
    \begin{tabular}{c|c|c}
      \multicolumn{3}{c}{\textbf{Table 10.1 Compact Groups}}\\
      group & rank & dim\\\hline
      $\mathrm{SU}(n+1)$ & $n$ & $n(n+2)$\\
      $\mathrm{SO}(2n+1)$ & $n$ & $n(2n+1)$\\
      $\mathrm{Sp}(n)$ & $n$ & $n(2n+1)$\\
      $\mathrm{SO}(2n)$ & $n$ & $n(2n-1)$
    \end{tabular}
    \qquad
    \begin{tabular}{c|c|c}
      \multicolumn{3}{c}{\textbf{Table 10.2 Noncompact Analogues}}\\
      (complexified group)/group & rank & dim\\\hline
      $\mathrm{SL}(n{+}1,\mathbb C)/\mathrm{SU}(n{+}1)$ & $n$ & $n(n+2)$\\
      $\mathrm{SO}(2n{+}1,\mathbb C)/\mathrm{SO}(2n{+}1)$ & $n$ & $n(2n+1)$\\
      $\mathrm{Sp}(n,\mathbb C)/\mathrm{Sp}(n)$ & $n$ & $n(2n+1)$\\
      $\mathrm{SO}(2n,\mathbb C)/\mathrm{SO}(2n)$ & $n$ & $n(2n-1)$
    \end{tabular}
  \end{center}
  \begin{center}
    \begin{tabular}{l|l|c|c|l}
      \multicolumn{5}{c}{\textbf{Table 10.3 Compact Homogeneous Spaces}}\\
      $\mathrm{Iso}$ & $\mathrm{Iso}_p$ & dim & rank & description\\\hline
      $\mathrm{SO}(n{+}1)$ & $\mathrm{SO}(n)$ & $n$ & $1$ & sphere\\
      $\mathrm O(n{+}1)$ & $\mathrm O(n)\times\{1,-1\}$ & $n$ & $1$ & $\mathbb{RP}^n$\\
      $\mathrm U(n{+}1)$ & $\mathrm U(n)\times\mathrm U(1)$ & $2n$ & $1$ & $\mathbb{CP}^n$\\
      $\mathrm{Sp}(n{+}1)$ & $\mathrm{Sp}(n)\times\mathrm{Sp}(1)$ & $4n$ & $1$ & $\mathbb{HP}^n$\\
      $\mathrm F_4$ & $\mathrm{Spin}(9)$ & $16$ & $1$ & $\mathbb{OP}^2$\\
      $\mathrm{SO}(p{+}q)$ & $\mathrm{SO}(p)\times\mathrm{SO}(q)$ & $pq$ & $\min(p,q)$ & real Grassmannian\\
      $\mathrm{SU}(p{+}q)$ & $\mathrm S(\mathrm U(p)\times\mathrm U(q))$ & $2pq$ & $\min(p,q)$ & complex Grassmannian\\
      $\mathrm{SU}(n)$ & $\mathrm{SO}(n)$ & $\tfrac{n^2-1}2$ & $n-1$ & $\mathbb R^n$'s in $\mathbb C^n$
    \end{tabular}
  \end{center}
  \begin{center}
    \begin{tabular}{l|l|c|c|l}
      \multicolumn{5}{c}{\textbf{Table 10.4 Noncompact Homogeneous Spaces}}\\
      $\mathrm{Iso}$ & $\mathrm{Iso}_p$ & dim & rank & description\\\hline
      $\mathrm{SO}(n,1)$ & $\mathrm{SO}(n)$ & $n$ & $1$ & hyperbolic space\\
      $\mathrm O(n,1)$ & $\mathrm O(n)\times\{1,-1\}$ & $n$ & $1$ & hyperbolic $\mathbb{RP}^n$\\
      $\mathrm U(n,1)$ & $\mathrm U(n)\times\mathrm U(1)$ & $2n$ & $1$ & hyperbolic $\mathbb{CP}^n$\\
      $\mathrm{Sp}(n,1)$ & $\mathrm{Sp}(n)\times\mathrm{Sp}(1)$ & $4n$ & $1$ & hyperbolic $\mathbb{HP}^n$\\
      $\mathrm F_4^{-20}$ & $\mathrm{Spin}(9)$ & $16$ & $1$ & hyperbolic $\mathbb{OP}^2$\\
      $\mathrm{SO}(p,q)$ & $\mathrm{SO}(p)\times\mathrm{SO}(q)$ & $pq$ & $\min(p,q)$ & hyperbolic Grassmannian\\
      $\mathrm{SU}(p,q)$ & $\mathrm S(\mathrm U(p)\times\mathrm U(q))$ & $2pq$ & $\min(p,q)$ & complex hyperbolic Grassmannian\\
      $\mathrm{SL}(n,\mathbb R)$ & $\mathrm{SO}(n)$ & $\tfrac{n^2-1}2$ & $n-1$ & Euclidean structures on $\mathbb R^n$
    \end{tabular}
  \end{center}
  Here $\mathrm{Spin}(n)$ is the universal double cover of $\mathrm{SO}(n)$,
  $n>2$, with the low-dimensional identities
  $\mathrm{SO}(2)=\mathrm U(1)$, $\mathrm{Spin}(3)=\mathrm{SU}(2)=\mathrm{Sp}(1)$,
  $\mathrm{Spin}(4)=\mathrm{Spin}(3)\times\mathrm{Spin}(3)$. By O'Neill's formula
  (Petersen's Theorem 4.5.3, outside this chapter) all the compact examples
  have $\sec\ge0$, and all the projective spaces (compact and noncompact) have
  quarter-pinched metrics, i.e.\ the ratio of smallest to largest sectional
  curvature is $\tfrac14$.
  \source{petersen:page-0379,petersen:page-0380}
\end{definition}

\begin{definition}[rank of a geodesic and of a Riemannian manifold; CROSSes]
  \label{def:pet-ch10-rank}
  \lean{PetersenLib.geodesicRank,PetersenLib.manifoldRank}
  The \emph{rank} of a geodesic $c:\mathbb R\to M$ is the dimension of the
  space of parallel fields $E$ along $c$ such that $R(E(t),\dot c(t))\dot c(t)=0$
  for all $t$; it is always $\ge1$, since $E=\dot c$ always qualifies. The
  \emph{rank} of a Riemannian manifold $(M,g)$ is the minimum rank over all
  geodesics of $M$. Any Cartesian product has rank $\ge2$, and many symmetric
  spaces have rank $\ge2$; all rank-one symmetric spaces are listed in
  \cref{def:pet-ch10-symmetric-space-tables} (Tables 10.3, 10.4), and the
  compact ones are called \emph{CROSSes}.
  \source{petersen:page-0381}
\end{definition}

\subsection{Isometries and Parallel Curvature}

Since a point symmetry $A_p$ preserves the curvature tensor and its covariant
derivative — $DA_p\big((\nabla_XR)(Y,Z,W)\big)=(\nabla_{DA_pX}R)(DA_pY,DA_pZ,DA_pW)$
— evaluating at $p$ (where $DA_p=-I$) gives
$-(\nabla_XR)(Y,Z,W)=(\nabla_{-X}R)(-Y,-Z,-W)=(\nabla_XR)(Y,Z,W)$, forcing
$\nabla R=0$ at $p$, and hence (by homogeneity, \cref{prop:pet-ch10-symmetric-homogeneous})
everywhere on $M$. This near-characterizes symmetric spaces.

\begin{definition}[locally symmetric space]
  \label{def:pet-ch10-locally-symmetric-space}
  \lean{PetersenLib.IsLocallySymmetric}
  A Riemannian manifold $(M,g)$ has \emph{parallel curvature tensor} if
  $\nabla R=0$; such a manifold is called \emph{locally symmetric}.
  \source{petersen:page-0382}
\end{definition}

\begin{theorem}[Thm. 10.1.1 (Cartan) — local symmetries from parallel curvature]
  \label{thm:pet-ch10-cartan-local-symmetry}
  \lean{PetersenLib.cartanLocalSymmetryExistence}
  \uses{def:pet-ch10-locally-symmetric-space,def:pet-ch10-symmetric-space}
  If $(M,g)$ is a locally symmetric Riemannian manifold, then for each $p\in M$
  there is an isometry $A_p$, defined on a neighborhood of $p$, with
  $DA_p=-I$ on $T_pM$. If moreover $(M,g)$ is simply connected and complete,
  the local symmetry is defined on all of $M$, and $(M,g)$ is a symmetric
  space.
  \source{petersen:page-0381,petersen:page-0382}
\end{theorem}
\begin{proof}
  \uses{thm:pet-ch10-cartan-local-symmetry,thm:pet-ch10-cartan-isometry-extension}
  \emph{Local statement.} Choose $\varepsilon>0$ small enough that
  $\exp_p:B(0,\varepsilon)\to B(p,\varepsilon)$ is a diffeomorphism, and define
  $A_p(x)=-x$ in these normal coordinates; the claim is that this candidate
  map is an isometry, i.e.\ that the metric in normal coordinates takes the
  same value at $x$ and $-x$. In polar coordinates on $B(0,\varepsilon)$ this
  reduces, by the fundamental equations relating curvature and the metric
  along radial geodesics, to showing that the curvature tensor $R$ agrees when
  transported the same distance in opposite radial directions from $p$. Two
  facts give this: first, at $p$ itself, $R(\cdot,v)v=R(\cdot,-v)(-v)$ for
  every $v\in T_pM$ (immediate from multilinearity and the symmetries of
  $R$), so the curvature tensors, expressed along the radial geodesics through
  $\pm v$, start out equal; second, since $\nabla R=0$, the radial field
  $\partial_r$ satisfies $\nabla_{\partial_r}R=0$ along every geodesic from
  $p$, so the curvature tensor along each radial geodesic satisfies the same
  (trivial) first-order ODE in $r$ with the same initial data at $r=0$ in the
  two opposite directions. Hence the curvature tensors agree at equal
  distances in opposite directions, giving equal metrics at $x$ and $-x$ and
  so an isometry $A_p$ near $p$.

  \emph{Global statement.} As in the proof of
  \cref{thm:pet-ch10-cartan-isometry-extension} below (and as in Petersen's
  Theorem 5.6.7, outside this chapter, for the analogous statement about
  constant-curvature isometries), the local isometry $A_p$ extends by
  analytic continuation along curves to a global isometry of $M$: completeness
  of $(M,g)$ guarantees the continuation never leaves the domain of
  definition, and simple connectivity of $M$ guarantees the continuation is
  independent of the path, hence single-valued. The resulting global isometry
  $A_p$ still satisfies $DA_p|_p=-I$, exhibiting $(M,g)$ as a symmetric space.
\end{proof}

\begin{remark}[left-invariant metrics need not be locally symmetric]
  \label{rem:pet-ch10-non-locally-symmetric-examples}
  \lean{PetersenLib.bergerSphereNotLocallySymmetric,PetersenLib.heisenbergNotLocallySymmetric}
  \uses{def:pet-ch10-locally-symmetric-space}
  The Berger spheres ($\varepsilon\ne1$) and the Heisenberg group, with their
  natural left-invariant metrics, do not have parallel curvature tensor and so
  are not locally symmetric; since both are $3$-dimensional, they cannot even
  have parallel Ricci tensor.
  \source{petersen:page-0382}
\end{remark}

\begin{theorem}[Thm. 10.1.2 (Cartan) — extending a curvature-preserving linear isometry]
  \label{thm:pet-ch10-cartan-isometry-extension}
  \lean{PetersenLib.cartanIsometryExtension}
  \uses{def:pet-ch10-symmetric-space,def:pet-ch10-locally-symmetric-space}
  Let $(M,g)$ be a simply connected symmetric space, $p\in M$, and let
  $(N,\bar g)$ be a complete locally symmetric space of the same dimension,
  $q\in N$. Given a linear isometry $L:T_pM\to T_qN$ such that
  $L(R^\sharp(x,y)z)=\bar R^\sharp(Lx,Ly)Lz$ for all $x,y,z\in T_pM$, there is
  a unique Riemannian isometry $F:M\to N$ with $D_pF=L$.
  \source{petersen:page-0382}
\end{theorem}
\begin{proof}
  \uses{thm:pet-ch10-cartan-isometry-extension}
  As in the constant-curvature case, the proof is by analytic continuation, so
  it suffices to construct $F$ locally near $p$. If a local isometry $F$ with
  $D_pF=L$ exists, it must be $F=\exp_q\circ L\circ\exp_p^{-1}$ on a normal
  neighborhood of $p$, since an isometry is determined by its value and
  differential at one point and this composite has the right value and
  differential at $p$. To check that this candidate $F$ is indeed an isometry,
  it suffices, exactly as in the proof of
  \cref{thm:pet-ch10-cartan-local-symmetry}, to show that the metrics of $M$
  and $N$ agree in exponential coordinates once the tangent spaces are
  identified by $L$; since radial curvatures determine the metric in such
  coordinates, it suffices to match the curvature tensors. Both $R$ and
  $\bar R$ are parallel (the first because $(M,g)$ is symmetric, by the
  computation preceding \cref{def:pet-ch10-locally-symmetric-space}, the
  second by hypothesis on $(N,\bar g)$), so both satisfy the same trivial
  first-order ODE $\nabla_{\partial_r}R=0$ along radial geodesics from $p$ and
  $q$ respectively; by hypothesis they agree initially, via $L$, at $p$ and
  $q$. Hence they remain equal (via $L$, extended by parallel translation)
  in radially parallel frames around $p$ and $q$, so the two metrics agree in
  exponential coordinates, and $F$ is a local isometry near $p$. The global
  statement and its uniqueness then follow by analytically continuing $F$
  along curves in $M$: simple connectivity of $M$ makes the continuation
  single-valued, and completeness of $N$ makes it everywhere defined.
\end{proof}

\subsection{The Lie Algebra Description}

Throughout this subsection $(M,g)$ is a symmetric space, $p\in M$, $A_p$ is the
point symmetry at $p$ (the isometric involution with $A_p(p)=p$, $DA_p|_p=-I$),
and $\mathrm{iso}=\mathrm{iso}(M,g)$ is the Lie algebra of Killing fields of $(M,g)$,
with isotropy subalgebra $\mathrm{iso}_p=\{X\in\mathrm{iso}\mid X|_p=0\}$. Since
$A_p$ is an isometry, its pushforward $(A_p)_*$ preserves Killing fields, giving a
natural linear map $\sigma_p=(A_p)_*:\mathrm{iso}\to\mathrm{iso}$,
$\sigma_p(X)=DA_p\circ X\circ A_p^{-1}$.

\begin{proposition}[Prop. 10.1.3 — the involution $\sigma_p$ on $\mathrm{iso}$]
  \label{prop:pet-ch10-sigma-p-involution}
  \lean{PetersenLib.sigmaPInvolution}
  The map $\sigma_p=(A_p)_*$ is an involution of $\mathrm{iso}$: $\sigma_p^2=\mathrm{id}$.
  Its $1$-eigenspace is $\mathrm{iso}_p$, and its $(-1)$-eigenspace is
  $\{X\in\mathrm{iso}\mid(\nabla X)|_p=0\}$.
  \source{petersen:page-0383,petersen:page-0384}
\end{proposition}
\begin{proof}
  \uses{prop:pet-ch10-sigma-p-involution}
  Since $A_p^2=\mathrm{id}$, clearly $\sigma_p^2=\mathrm{id}$, so $\mathrm{iso}$
  is the direct sum of the $(\pm1)$-eigenspaces of $\sigma_p$.

  Because $A_p$ is an isometry, pushforward commutes with the Levi-Civita
  connection: $(A_p)_*(\nabla_VX)=\nabla_{(A_p)_*V}(A_p)_*X$ for all vector fields
  $V,X$. Fix $v\in T_pM$ and choose $V$ with $V|_p=-v$; since $A_p(p)=p$ and
  $DA_p|_p=-I$, $(A_p)_*V|_p=DA_p(V|_p)=-V|_p=v$, so evaluating the identity
  above at $p$ gives, for every $X\in\mathrm{iso}$,
  \[
    (\nabla_v\sigma_p(X))|_p=DA_p\big((\nabla_VX)|_p\big)=-(\nabla_{-v}X)|_p=(\nabla_vX)|_p,
  \]
  i.e.\ $(\nabla\sigma_p(X))|_p=(\nabla X)|_p$ for every $X\in\mathrm{iso}$
  \emph{(unconditionally, not just on an eigenspace)}. Also, evaluating
  $\sigma_p(X)=DA_p\circ X\circ A_p^{-1}$ at $p$ gives $\sigma_p(X)|_p=DA_p(X|_p)=-X|_p$.

  \emph{$1$-eigenspace.} If $\sigma_p(X)=X$ then $X|_p=\sigma_p(X)|_p=-X|_p$, so
  $X|_p=0$, i.e.\ $X\in\mathrm{iso}_p$. Conversely, if $X\in\mathrm{iso}_p$ then
  $\sigma_p(X)|_p=-X|_p=0$ too, so $\sigma_p(X)\in\mathrm{iso}_p$ and $X,\sigma_p(X)$
  agree at $p$; by the identity above they also have the same covariant derivative
  at $p$. A Killing field is determined by its value and covariant derivative at one
  point, so $\sigma_p(X)=X$.

  \emph{$(-1)$-eigenspace.} If $\sigma_p(X)=-X$ then, using the unconditional
  identity, $(\nabla X)|_p=(\nabla\sigma_p(X))|_p=(\nabla(-X))|_p=-(\nabla X)|_p$,
  forcing $(\nabla X)|_p=0$. Conversely, if $(\nabla X)|_p=0$, then $-X$ and
  $\sigma_p(X)$ agree at $p$ (both equal $-X|_p$) and both have vanishing covariant
  derivative there (for $-X$ this is immediate, and for $\sigma_p(X)$ it follows from
  the unconditional identity applied to $X$), so $\sigma_p(X)=-X$ by the same
  uniqueness principle.
\end{proof}

Petersen's Proposition 8.1.4 (not among the transcribed pages) supplies the short
exact sequence $0\to\mathrm{iso}_p\to\mathrm{iso}\to t_p\to0$, $t_p=\{X|_p\in T_pM\mid
X\in\mathrm{iso}\}$; since $M$ is homogeneous, $t_p=T_pM$, and since $M$ is symmetric
\cref{prop:pet-ch10-sigma-p-involution} shows the $(-1)$-eigenspace of $\sigma_p$
maps isomorphically onto $T_pM$ under evaluation at $p$.

\begin{definition}[the subspaces $t_p$ and $s_p$]
  \label{def:pet-ch10-tp-decomposition}
  \lean{PetersenLib.tpSubspace,PetersenLib.spSubspace}
  \uses{prop:pet-ch10-sigma-p-involution}
  Redefine $t_p\subset\mathrm{iso}$ as the $(-1)$-eigenspace of $\sigma_p$, i.e.\
  $t_p=\{X\in\mathrm{iso}\mid(\nabla X)|_p=0\}$; by
  \cref{prop:pet-ch10-sigma-p-involution} evaluation at $p$ is a linear isomorphism
  $t_p\xrightarrow{\sim}T_pM$. This gives the natural decomposition
  $\mathrm{iso}=t_p\oplus\mathrm{iso}_p$, and evaluating at $p$ the alternate
  representation
  \[
    \mathrm{iso}\simeq T_pM\oplus s_p\subset T_pM\oplus\mathfrak{so}(T_pM),
    \qquad
    s_p=\{(\nabla X)|_p\in\mathfrak{so}(T_pM)\mid X\in\mathrm{iso}_p\}.
  \]
  \source{petersen:page-0384}
\end{definition}

\begin{proposition}[Prop. 10.1.4 — the bracket of $\mathrm{iso}$ from the curvature]
  \label{prop:pet-ch10-lie-bracket-curvature}
  \lean{PetersenLib.symmetricSpaceLieBracketCurvature}
  \uses{def:pet-ch10-tp-decomposition}
  Under the identification $\mathrm{iso}\simeq T_pM\oplus s_p$ of
  \cref{def:pet-ch10-tp-decomposition}, the Lie algebra structure of $\mathrm{iso}$
  is determined entirely by the curvature tensor at $p$:
  \begin{enumerate}
    \item If $X,Y\in T_pM$, then $[X,Y]=R(X,Y)\in s_p$.
    \item If $X\in T_pM$ and $S\in s_p$, then $[X,S]=-S(X)\in T_pM$.
    \item If $S,T\in s_p$, then $[S,T]=-(S\circ T-T\circ S)\in s_p$.
  \end{enumerate}
  \source{petersen:page-0384,petersen:page-0385}
\end{proposition}
\begin{proof}
  \uses{prop:pet-ch10-lie-bracket-curvature}
  We rely on the Killing-field curvature identity (Petersen's Proposition 8.1.3):
  for $X,Y\in\mathrm{iso}$, $[\nabla X,\nabla Y](V)+\nabla_V[X,Y]=R(X,Y)V$ for all
  vector fields $V$.

  \emph{(1)} For a Killing field, the torsion-free identity $[X,Y]=\nabla_XY-\nabla_YX$
  holds; if $X,Y\in t_p$ then $(\nabla X)|_p=(\nabla Y)|_p=0$, so $[X,Y]|_p=0$,
  i.e.\ $[X,Y]\in\mathrm{iso}_p$. Apply the curvature identity at $p$ with these
  $X,Y$: since $(\nabla X)|_p=(\nabla Y)|_p=0$, also $[\nabla X,\nabla Y]|_p=0$,
  giving $(\nabla[X,Y])|_p(Z|_p)=R(X|_p,Y|_p)Z|_p$ for all $Z\in T_pM$. Thus
  $(\nabla[X,Y])|_p=R(X|_p,Y|_p)\in\mathfrak{so}(T_pM)$, i.e.\ $[X,Y]=R(X,Y)\in s_p$
  under the identification.

  \emph{(2)} For $X\in t_p$, $Y\in\mathrm{iso}_p$ (so $Y|_p=0$ and
  $(\nabla Y)|_p=S\in s_p$), the torsion-free identity gives $[X,Y]=\nabla_YX-\nabla_XY$,
  and at $p$: $(\nabla_YX)|_p=0$ since $Y|_p=0$, while $(\nabla_XY)|_p=S(X|_p)$.
  Hence $[X,Y]|_p=-S(X|_p)$, i.e.\ $[X,S]=-S(X)$.

  \emph{(3)} For $X,Y\in\mathrm{iso}_p$ (so $X|_p=Y|_p=0$, with
  $(\nabla X)|_p=S$, $(\nabla Y)|_p=T$), $[X,Y]=\nabla_XY-\nabla_YX$ vanishes at
  $p$, so $[X,Y]\in\mathrm{iso}_p$ again. Applying the curvature identity with
  $V=Z\in T_pM$ at $p$: since $X|_p=Y|_p=0$, $R(X|_p,Y|_p)Z|_p=0$, so
  $[\nabla X,\nabla Y](Z)|_p+(\nabla[X,Y])|_p(Z)=0$, i.e.\
  $(\nabla[X,Y])|_p=-[\nabla X,\nabla Y]|_p=-[S,T]=-(S\circ T-T\circ S)$.
\end{proof}

\begin{definition}[adjoint action and Killing form]
  \label{def:pet-ch10-killing-form}
  \lean{PetersenLib.adjointAction,PetersenLib.killingForm}
  On a (finite-dimensional) Lie algebra $\mathfrak g$ the \emph{adjoint action} is
  $\mathrm{ad}_X:\mathfrak g\to\mathfrak g$, $\mathrm{ad}_X(Y)=[X,Y]$, and the
  \emph{Killing form} is $B(X,Y)=\operatorname{tr}(\mathrm{ad}_X\circ\mathrm{ad}_Y)$.
  The Killing form is symmetric, and $\mathrm{ad}_X$ is skew-symmetric with
  respect to $B$: $B(\mathrm{ad}_X(Y),Z)+B(Y,\mathrm{ad}_X(Z))=0$.
  \source{petersen:page-0385}
\end{definition}

\begin{remark}[Rem. 10.1.5 — $\mathrm{ad}_S$ is skew and $B\le0$ on $s_p$]
  \label{rem:pet-ch10-adjoint-skew-symmetric}
  \lean{PetersenLib.adjointOnSpNegativeSemidefinite}
  \uses{prop:pet-ch10-lie-bracket-curvature,def:pet-ch10-killing-form}
  Equip $\mathrm{iso}\simeq T_pM\oplus s_p$ with the natural inner product
  $g((X_1,S_1),(X_2,S_2))=g(X_1,X_2)+g(S_1,S_2)$, where $g(S_1,S_2)=-\operatorname{tr}(S_1\circ S_2)$
  is the standard inner product on $\mathfrak{so}(T_pM)$. For $S\in s_p$ the map
  $\mathrm{ad}_S:\mathrm{iso}\to\mathrm{iso}$ is skew-symmetric with respect to this
  inner product, so its adjoint satisfies $(\mathrm{ad}_S)^*=-\mathrm{ad}_S$; hence
  \[
    B(S,S)=\operatorname{tr}(\mathrm{ad}_S\circ\mathrm{ad}_S)=-\operatorname{tr}\big(\mathrm{ad}_S\circ(\mathrm{ad}_S)^*\big)\le0,
  \]
  with equality only if $S=0$.
  \source{petersen:page-0385}
\end{remark}
\begin{proof}
  \uses{rem:pet-ch10-adjoint-skew-symmetric}
  The inequality $B(S,S)=-\operatorname{tr}(\mathrm{ad}_S\circ(\mathrm{ad}_S)^*)\le0$ holds
  because $\operatorname{tr}(A\circ A^*)\ge0$ for any linear map $A$ on an inner-product
  space (it is the squared Hilbert–Schmidt norm of $A$). If $B(S,S)=0$ then
  $\operatorname{tr}(\mathrm{ad}_S\circ(\mathrm{ad}_S)^*)=0$ forces $\mathrm{ad}_S=0$. In
  particular $\mathrm{ad}_S(X)=[S,X]=S(X)=0$ for every $X\in T_pM$ (by
  \cref{prop:pet-ch10-lie-bracket-curvature}(2)), and since $S\in\mathfrak{so}(T_pM)$
  is itself a linear map on $T_pM$, vanishing on all of $T_pM$ forces $S=0$.
\end{proof}

\begin{theorem}[Thm. 10.1.6 — curvature from the bracket]
  \label{thm:pet-ch10-curvature-bracket-formula}
  \lean{PetersenLib.curvatureFromBracket,PetersenLib.ricciFromKillingForm}
  \uses{prop:pet-ch10-lie-bracket-curvature,def:pet-ch10-killing-form}
  If $X,Y,Z\in t_p$, then at $p$, $R(X,Y)Z=[Z,[X,Y]]$; equivalently, under
  $\mathrm{iso}\simeq T_pM\oplus s_p$,
  \[
    R(X,Y)Z=-\mathrm{ad}_Z\circ\mathrm{ad}_Y(X),
    \qquad
    \operatorname{Ric}(Y,Z)=-\tfrac12B(Z,Y).
  \]
  Moreover, if $\operatorname{Ric}=\lambda g$ for a scalar $\lambda\ne0$, the curvature
  operator has the same sign as $\lambda$.
  \source{petersen:page-0385,petersen:page-0386,petersen:page-0387}
\end{theorem}
\begin{proof}
  \uses{thm:pet-ch10-curvature-bracket-formula,rem:pet-ch10-adjoint-skew-symmetric}
  We use the identity $\nabla^2_{A,B}C=-R(C,A)B$ valid for Killing fields
  $A,B,C\in\mathrm{iso}$ (Petersen's Proposition 8.1.3). For $X,Y,Z\in t_p$,
  $(\nabla X)|_p=(\nabla Y)|_p=(\nabla Z)|_p=0$. Applying the identity with
  $(A,B,C)=(Z,X,Y)$ and $(A,B,C)=(Z,Y,X)$, and using that the connection terms
  vanish at $p$ (as $(\nabla Z)|_p=0$), gives at $p$
  \[
    \nabla_Z\nabla_XY=-R(Y,Z)X,\qquad\nabla_Z\nabla_YX=-R(X,Z)Y.
  \]
  Subtracting and using $[X,Y]=\nabla_XY-\nabla_YX$, together with the first
  Bianchi identity $R(X,Y)Z+R(Y,Z)X+R(Z,X)Y=0$, gives at $p$
  \[
    \nabla_Z[X,Y]=-R(Y,Z)X+R(X,Z)Y=R(X,Y)Z=[Z,[X,Y]],
  \]
  the last equality by the identification of $\nabla_Z[X,Y]|_p$ with the bracket of
  the Killing fields $Z\in t_p$ and $[X,Y]\in\mathrm{iso}_p$ furnished by
  \cref{prop:pet-ch10-lie-bracket-curvature}. Writing $[X,Y]=\mathrm{ad}_Y(X)\in s_p$
  under the identification and applying \cref{prop:pet-ch10-lie-bracket-curvature}(2)
  once more gives $R(X,Y)Z=[Z,\mathrm{ad}_Y(X)]=-\mathrm{ad}_Z\circ\mathrm{ad}_Y(X)$.

  For the Ricci formula, note $\mathrm{ad}_Z\circ\mathrm{ad}_Y$ leaves the
  decomposition $T_pM\oplus s_p$ invariant (each of $\mathrm{ad}_Z,\mathrm{ad}_Y$
  interchanges the two summands, by parts (1)–(2) of
  \cref{prop:pet-ch10-lie-bracket-curvature}), so
  \[
    \operatorname{tr}\big(\mathrm{ad}_Z\circ\mathrm{ad}_Y|_{T_pM}\big)
    =\operatorname{tr}\big(\mathrm{ad}_Z|_{s_p}\circ\mathrm{ad}_Y|_{T_pM}\big)
    =\operatorname{tr}\big(\mathrm{ad}_Y|_{T_pM}\circ\mathrm{ad}_Z|_{s_p}\big)
    =\operatorname{tr}\big(\mathrm{ad}_Y\circ\mathrm{ad}_Z|_{s_p}\big),
  \]
  using cyclicity of the trace. By symmetry of $B$,
  \[
    B(Z,Y)=\operatorname{tr}(\mathrm{ad}_Z\circ\mathrm{ad}_Y)
    =\operatorname{tr}(\mathrm{ad}_Z\circ\mathrm{ad}_Y|_{T_pM})+\operatorname{tr}(\mathrm{ad}_Z\circ\mathrm{ad}_Y|_{s_p})
    =2\operatorname{tr}(\mathrm{ad}_Z\circ\mathrm{ad}_Y|_{T_pM}),
  \]
  and the curvature formula above shows
  $\operatorname{Ric}(Z,Y)=-\operatorname{tr}(\mathrm{ad}_Z\circ\mathrm{ad}_Y|_{T_pM})=-\tfrac12B(Z,Y)$.

  Finally, suppose $\operatorname{Ric}=\lambda g$, $\lambda\ne0$. Using
  $\operatorname{Ric}(Y,Z)=-\tfrac12B(Z,Y)$ to rewrite the Killing form in terms of
  $\operatorname{Ric}$, and then in terms of $g$, and the bracket identities of
  \cref{prop:pet-ch10-lie-bracket-curvature} to move the $\operatorname{ad}$'s across the
  pairing, one computes for $X,Y,W,Z\in T_pM$
  \[
    g(R(X,Y)Z,W)|_p
    =-\tfrac1{2\lambda}B(\mathrm{ad}_Y(X),\mathrm{ad}_Z(W))
    =-\tfrac1{2\lambda}B([X,Y],[W,Z])
    =\tfrac1{2\lambda}\operatorname{tr}\big(\mathrm{ad}_{[X,Y]}\circ(\mathrm{ad}_{[W,Z]})^*\big).
  \]
  Summing over an orthonormal basis of a $2$-plane, $X_i\wedge Y_i$, the diagonal
  term of the curvature operator becomes
  \[
    g\big(\mathfrak R(\textstyle\sum X_i\wedge Y_i),\sum X_i\wedge Y_i\big)\big|_p
    =\sum g(R(X_i,Y_i)Y_j,X_j)|_p
    =-\tfrac1{2\lambda}B\Big(\sum\mathrm{ad}_{X_i}(Y_i),\sum\mathrm{ad}_{X_i}(Y_i)\Big).
  \]
  Since $\sum\mathrm{ad}_{X_i}(Y_i)\in s_p$, \cref{rem:pet-ch10-adjoint-skew-symmetric}
  gives $B(\sum\mathrm{ad}_{X_i}(Y_i),\sum\mathrm{ad}_{X_i}(Y_i))\le0$, so every
  diagonal term of the curvature operator has the same sign as $\lambda$; hence so
  does the curvature operator itself.
\end{proof}

\begin{definition}[the abstract bracket on $T_pM\oplus\mathfrak{so}(T_pM)$]
  \label{def:pet-ch10-abstract-symmetric-bracket}
  \lean{PetersenLib.abstractSymmetricSpaceBracket}
  \uses{thm:pet-ch10-curvature-bracket-formula}
  For a Riemannian manifold $(M,g)$ (not assumed symmetric) and $p\in M$, define
  a bracket on $\mathfrak N_p:=T_pM\oplus\mathfrak{so}(T_pM)$ by
  \[
    [X,Y]=R(X,Y)\in\mathfrak{so}(T_pM)\quad(X,Y\in T_pM),
  \]
  \[
    -[S,X]=[X,S]=S(X)\in T_pM\quad(X\in T_pM,\ S\in\mathfrak{so}(T_pM)),
  \]
  \[
    [S,S']=-(S\circ S'-S'\circ S)\in\mathfrak{so}(T_pM)\quad(S,S'\in\mathfrak{so}(T_pM)).
  \]
  This bracket need not satisfy the Jacobi identity on triples involving exactly two
  elements of $T_pM$, where the Jacobi expression evaluates to
  $-S\circ R_{X,Y}+R_{X,Y}\circ S+R_{S(X),Y}+R_{X,S(Y)}$ for $X,Y\in T_pM$,
  $S\in\mathfrak{so}(T_pM)$; it does hold on triples with zero or one elements from
  $T_pM$ (immediate from skew-symmetry of the bracket on
  $\mathfrak{so}(T_pM)$), and on triples of three elements of $T_pM$ it is
  equivalent to the first Bianchi identity
  $R(X,Y)Z+R(Z,X)Y+R(Y,Z)X=0$, i.e.\ $[Z,[X,Y]]+[Y,[Z,X]]+[X,[Y,Z]]=0$.
  \source{petersen:page-0387}
\end{definition}

\begin{corollary}[Cor. 10.1.7 — membership criterion for $s_p$]
  \label{cor:pet-ch10-so-membership-criterion}
  \lean{PetersenLib.spMembershipCriterion}
  \uses{def:pet-ch10-tp-decomposition}
  Let $(M,g)$ be a simply connected symmetric space. Then $S\in\mathfrak{so}(T_pM)$
  lies in $s_p$ if and only if, for all $X,Y\in T_pM$,
  \[
    -S\circ R_{X,Y}+R_{X,Y}\circ S+R_{S(X),Y}+R_{X,S(Y)}=0.
  \]
  \source{petersen:page-0388}
\end{corollary}
\begin{proof}
  \uses{cor:pet-ch10-so-membership-criterion,thm:pet-ch10-cartan-isometry-extension}
  ($\Rightarrow$) If $Z\in\mathrm{iso}_p$ and $S=(\nabla Z)|_p$, then $Z$ Killing
  gives $\mathcal L_ZR=0$; since $[V,Z]|_p=S(V)$ for $V\in T_pM$, evaluating
  $(\mathcal L_ZR)_{X,Y}V=0$ at $p$ gives
  \[
    0=-S(R_{X,Y}V)+R_{X,Y}S(V)+R_{S(X),Y}V+R_{X,S(Y)}V,
  \]
  for all $V\in T_pM$, which is the stated identity.

  ($\Leftarrow$) Given $S\in\mathfrak{so}(T_pM)$ satisfying the identity, construct
  the one-parameter isometric flow of $S$ on $T_pM$.
  \cref{thm:pet-ch10-cartan-isometry-extension} (Theorem 10.1.2) shows that a
  curve of linear isometries of $T_pM$ integrates to a curve of global isometries
  of $M$ exactly when the corresponding infinitesimal condition on the curvature
  tensor holds; the hypothesis on $S$ is exactly that condition, so the flow of
  $S$ generates a one-parameter group of isometries of $M$, whose infinitesimal
  generator is a Killing field $Z\in\mathrm{iso}_p$ with $(\nabla Z)|_p=S$.
\end{proof}

\begin{corollary}[Cor. 10.1.8 — $T_pM\oplus s_p$ is the maximal Lie subalgebra]
  \label{cor:pet-ch10-maximal-subalgebra}
  \lean{PetersenLib.tpSpMaximalSubalgebra}
  \uses{def:pet-ch10-abstract-symmetric-bracket,cor:pet-ch10-so-membership-criterion}
  Let $(M,g)$ be a simply connected symmetric space. The bracket of
  \cref{def:pet-ch10-abstract-symmetric-bracket} makes $T_pM\oplus s_p$ into a Lie
  algebra, with subalgebra $c_p:=T_pM\oplus\{0\}$. In fact $T_pM\oplus s_p$ is the
  maximal such Lie algebra: $c_p\subset T_pM\oplus s_p\subset\mathfrak N_{\mathfrak N_p}(c_p)$,
  where $\mathfrak N_{\mathfrak N_p}(c_p)$ denotes the normalizer of $c_p$ in
  $\mathfrak N_p=T_pM\oplus\mathfrak{so}(T_pM)$. The Lie algebra involution on
  $T_pM\oplus s_p$ (and on $c_p$) has $T_pM$ as $(-1)$-eigenspace and $s_p$ as
  $1$-eigenspace.
  \source{petersen:page-0388}
\end{corollary}
\begin{proof}
  \uses{cor:pet-ch10-maximal-subalgebra}
  By \cref{cor:pet-ch10-so-membership-criterion}, any subalgebra
  $\mathfrak t\subset\mathfrak{so}(T_pM)$ such that $T_pM\oplus\mathfrak t$ is a Lie
  subalgebra of $\mathfrak N_p$ (i.e.\ also satisfies the Jacobi identity on the
  mixed triples of \cref{def:pet-ch10-abstract-symmetric-bracket}) must consist of
  elements $S$ satisfying the criterion, hence $\mathfrak t\subset s_p$; this gives
  $T_pM\oplus\mathfrak t\subset T_pM\oplus s_p$, so $T_pM\oplus s_p$ is maximal
  among such subalgebras. On the other hand $T_pM\oplus s_p\subset\mathfrak N_{\mathfrak N_p}(c_p)$
  because for $S\in s_p$ and $X\in T_pM$, $[S,X]=-S(X)\in T_pM=c_p$ and
  $[X,X']=R(X,X')\in s_p$ normalizes $c_p$ by the same computation.
\end{proof}

\begin{remark}[reconstructing symmetric spaces from an involutive Lie algebra]
  \label{rem:pet-ch10-symmetric-space-from-involution}
  \lean{PetersenLib.symmetricSpaceFromInvolutiveLieAlgebra}
  \uses{cor:pet-ch10-maximal-subalgebra}
  Conversely, let $\mathfrak g$ be a Lie algebra with a Lie algebra involution
  $\sigma:\mathfrak g\to\mathfrak g$, and decompose $\mathfrak g=\mathfrak t\oplus\mathfrak k$
  into the $(-1)$- and $1$-eigenspaces of $\sigma$. Then $\mathfrak k$ is a Lie
  subalgebra (since $\sigma[X,Y]=[\sigma X,\sigma Y]=[X,Y]$ for $X,Y\in\mathfrak k$),
  and $[\mathfrak k,\mathfrak t]\subset\mathfrak t$, $[\mathfrak t,\mathfrak t]\subset\mathfrak k$.
  Suppose there is a connected compact Lie group $K$ with Lie algebra $\mathfrak k$
  whose bracket action on $\mathfrak t$ integrates to an action of $K$ on
  $\mathfrak t$ (automatic if $K$ is simply connected); compactness lets one choose
  a $K$-invariant Euclidean metric on $\mathfrak t$, making
  $\mathfrak g=\mathfrak t\oplus\mathfrak k$ of the type $\mathrm{iso}=t_p\oplus\mathrm{iso}_p$.
  If further there is a Lie group $G\supset K$ with Lie algebra $\mathfrak g$, one
  obtains a Riemannian manifold $G/K$; when $G$ is simply connected, $\sigma$ is
  the differential of a Lie group involution $A:G\to G$ fixing $K$ pointwise,
  giving the point symmetry at $p=eK$ that makes $G/K$ a symmetric space.
  Without connectedness or simple-connectivity hypotheses there is a long exact
  sequence of homotopy groups
  \[
    \pi_1(K)\to\pi_1(G)\to\pi_1(G/K)\to\pi_0(K)\to\pi_0(G)\to\pi_0(G/K)\to1,
  \]
  from which $G/K$ is connected and simply connected once $\pi_0(K)\to\pi_0(G)$
  is an isomorphism and $\pi_1(K)\to\pi_1(G)$ is surjective. The involution
  $\sigma$ is an essential part of the data: it need not exist on a given Lie
  algebra, and $-\mathrm{id}$ is never a valid choice on a nonabelian $\mathfrak g$,
  since it is a Lie algebra \emph{anti}-automorphism, not an automorphism.
  \source{petersen:page-0388,petersen:page-0389}
\end{remark}

\section{Examples of Symmetric Spaces}

Throughout this section, for $X,Y\in\operatorname{Mat}_{k\times l}(\mathbb F)$ ($\mathbb F=\mathbb R$
or $\mathbb C$) write $\langle X,Y\rangle=\operatorname{tr}(X^\top Y)=\overline{\operatorname{tr}(XY^*)}$
(conjugation relevant only for $\mathbb F=\mathbb C$); this inner product is invariant
under $(A,B,X)\mapsto AXB^{-1}=AXB^\top$ for $A\in O(k)$, $B\in O(l)$, since
$\langle O_1XO_2,O_1YO_2\rangle=\operatorname{tr}(O_2^\top X^\top O_1^\top O_1YO_2)=\operatorname{tr}(X^\top Y)$.

\begin{example}[§10.2.1 — the compact Grassmannian]
  \label{ex:pet-ch10-compact-grassmannian}
  \lean{PetersenLib.compactGrassmannianSymmetricSpace}
  \uses{def:pet-ch10-tp-decomposition,def:pet-ch10-killing-form,thm:pet-ch10-curvature-bracket-formula,cor:pet-ch10-maximal-subalgebra}
  For $k,l\ge1$ with $k+l\ge3$, split $\mathbb R^{k+l}=\mathbb R^k\oplus\mathbb R^l$ and let
  $M=\widetilde G_k(\mathbb R^{k+l})$ be the Grassmannian of \emph{oriented}
  $k$-planes in $\mathbb R^{k+l}$, with distinguished point $p=\mathbb R^k$ (positively
  oriented). Then $M$ is the homogeneous space $O(k+l)/(\mathrm{SO}(k)\times O(l))$,
  $T_pM\cong\operatorname{Mat}_{k\times l}(\mathbb R)$ (an $\mathbb R^k$-graph over the
  complementary $\mathbb R^l$), the isotropy action of $\mathrm{SO}(k)\times O(l)$ on
  $\operatorname{Mat}_{k\times l}$ is $(A,B,X)\mapsto AXB^\top$, and $M$ is a symmetric
  space with involution induced by $(A,B)=(I_k,-I_l)$. In the Lie algebra
  decomposition $\mathfrak{so}(k+l)=t_p\oplus\mathfrak{so}(k)\oplus\mathfrak{so}(l)$,
  with
  \[
    t_p=\Big\{\begin{pmatrix}0&B\\-B^\top&0\end{pmatrix}\mid B\in\operatorname{Mat}_{k\times l}\Big\},
  \]
  the natural inner product $\langle X,Y\rangle=\operatorname{tr}(X^\top Y)$ on $t_p$ (skew
  matrices, $=-\operatorname{tr}(XY)$) equals $-\tfrac1{k+l-2}B(X,Y)$, and
  \[
    \operatorname{Ric}=\tfrac{k+l-2}2g,\qquad
    \langle R(X,Y)Y,X\rangle=|[X,Y]|^2\ge0.
  \]
  Moreover $t_p=\mathfrak{so}(k)\otimes\mathfrak{so}(l)$-irreducible (but the isotropy
  action, though irreducible, is not transitive on the unit sphere unless $k=1$ or
  $l=1$, in which case the metric has constant positive curvature), and
  \cref{cor:pet-ch10-maximal-subalgebra} identifies $\mathrm{iso}=\mathfrak{so}(k+l)$.
  \source{petersen:page-0390,petersen:page-0391,petersen:page-0392}
\end{example}
\begin{proof}
  \uses{ex:pet-ch10-compact-grassmannian}
  \emph{Homogeneous structure.} $O(k+l)$ acts transitively on positively oriented
  orthonormal $k$-frames, hence on oriented $k$-planes, by $O\cdot(e_1,\dots,e_k)=
  (Oe_1,\dots,Oe_k)$; the isotropy at $p=\mathbb R^k$ is exactly $\mathrm{SO}(k)\times O(l)$
  (orientation-preserving on $\mathbb R^k$, arbitrary orthogonal on $\mathbb R^l$).
  A $k$-plane near $p$ is the graph of a linear map $\mathbb R^k\to\mathbb R^l$, giving
  $T_pM\cong\operatorname{Mat}_{k\times l}$; conjugating this graph representation by
  $(A,B)\in\mathrm{SO}(k)\times O(l)$ replaces the graph map $X$ by $AXB^{-1}=AXB^\top$.

  \emph{Involution.} Taking $A=I_k$, $B=-I_l$ gives $AXB^\top=-X$ for all $X$,
  so $(I_k,-I_l)\in\mathrm{SO}(k)\times O(l)$ acts as $-1$ on $T_pM$, exhibiting $M$
  as a symmetric space (without yet identifying $\mathrm{iso}$).

  \emph{Lie algebra decomposition.} Write $X\in\mathfrak{so}(k+l)$ in blocks
  $X=\begin{pmatrix}X_1&B\\-B^\top&X_2\end{pmatrix}$ with $X_1\in\mathfrak{so}(k)$,
  $X_2\in\mathfrak{so}(l)$, $B\in\operatorname{Mat}_{k\times l}$; this is exactly the
  orthogonal decomposition $\mathfrak{so}(k+l)=t_p\oplus\mathfrak{so}(k)\oplus\mathfrak{so}(l)$
  with $t_p=T_pM$ the off-diagonal block, on which the involution above acts as
  $-1$.

  \emph{Metric and curvature.} Since $t_p$ consists of skew-symmetric matrices,
  $\langle X,Y\rangle=\operatorname{tr}(X^\top Y)=-\operatorname{tr}(XY)$; on the other hand the
  Killing form of $\mathfrak{so}(n)$ is $B(X,Y)=(n-2)\operatorname{tr}(XY)$ (Exercise
  10.5.6(5)), so with $n=k+l$, $\langle X,Y\rangle=-\tfrac1{k+l-2}B(X,Y)$. Combined
  with \cref{thm:pet-ch10-curvature-bracket-formula}
  ($\operatorname{Ric}(Y,Z)=-\tfrac12B(Z,Y)$), this gives $\operatorname{Ric}=\tfrac{k+l-2}2g$, and
  \[
    \langle R(X,Y)Y,X\rangle=\langle[X,Y],[X,Y]\rangle=-\tfrac1{k+l-2}B([X,Y],[X,Y])=|[X,Y]|^2\ge0
  \]
  by \cref{thm:pet-ch10-curvature-bracket-formula} again (curvature has the sign
  of the Einstein constant, here positive).

  \emph{Isometry group.} Since the isotropy of $t_p$ is irreducible, any subalgebra
  properly containing $\mathfrak{so}(k)\oplus\mathfrak{so}(l)$ inside $\mathfrak{so}(k+l)$
  and normalizing it must be all of $\mathfrak{so}(k+l)$; \cref{cor:pet-ch10-maximal-subalgebra}
  then gives $\mathrm{iso}=\mathfrak{so}(k+l)$, so $\mathrm{SO}(k+l)$ is a covering of
  $\mathrm{Iso}(M,g)$.
\end{proof}

\begin{example}[§10.2.2 — the hyperbolic Grassmannian]
  \label{ex:pet-ch10-hyperbolic-grassmannian}
  \lean{PetersenLib.hyperbolicGrassmannianSymmetricSpace}
  \uses{ex:pet-ch10-compact-grassmannian}
  For $k,l>0$ let $\mathbb R^{k,l}$ carry the quadratic form
  $v^\top I_{k,l}w=\sum_{i\le k}v_iw_i-\sum_{i>k}v_iw_i$, with isometry group
  $O(k,l)=\{X\mid X\,I_{k,l}\,X^\top=I_{k,l}\}$ (noncompact when $k,l>0$,
  containing the maximal compact subgroup $O(k)\times O(l)$) and Lie algebra
  $\mathfrak{so}(k,l)=\{Y\mid Y\,I_{k,l}+I_{k,l}\,Y^\top=0\}=\big\{\begin{pmatrix}Y_1&B\\B^\top&Y_2\end{pmatrix}
  \mid Y_1\in\mathfrak{so}(k),\,Y_2\in\mathfrak{so}(l),\,B\in\operatorname{Mat}_{k\times l}\big\}$.
  The \emph{hyperbolic Grassmannian} $M=\widetilde G_k(\mathbb R^{k,l})$ is the space of
  oriented $k$-planes on which the form is positive definite (an open subset of
  $\widetilde G_k(\mathbb R^{k+l})$); $O(k,l)$ acts transitively with isotropy
  $\mathrm{SO}(k)\times O(l)$ at $p=\mathbb R^k$, and $M$ is symmetric with
  $\mathfrak{so}(k,l)=t_p\oplus\mathfrak{so}(k)\oplus\mathfrak{so}(l)$,
  $t_p=\big\{\begin{pmatrix}0&B\\B^\top&0\end{pmatrix}\big\}$ now consisting of
  \emph{symmetric} matrices, $\langle X,Y\rangle=\operatorname{tr}(XY)=\tfrac1{k+l-2}B(X,Y)$,
  \[
    \operatorname{Ric}=-\tfrac{k+l-2}2g,\qquad
    \langle R(X,Y)Y,X\rangle=-|[X,Y]|^2\le0,
  \]
  and $\mathrm{iso}=\mathfrak{so}(k,l)$. When $l=1$ this recovers hyperbolic space
  $H^k$ with its full isometry group.
  \source{petersen:page-0392,petersen:page-0393}
\end{example}
\begin{proof}
  \uses{ex:pet-ch10-hyperbolic-grassmannian}
  The homogeneous and symmetric-space structure is identical to
  \cref{ex:pet-ch10-compact-grassmannian}, replacing $O(k+l)$ by $O(k,l)$
  throughout, since $O(k,l)\supset O(k)\times O(l)$ acts transitively on
  positive-definite oriented $k$-planes and the same involution $(I_k,-I_l)$ fixes
  $p=\mathbb R^k$ and acts as $-1$ on $T_pM$. The only change is that $t_p$ now
  consists of matrices \emph{symmetric} with respect to transpose (from the $+B^\top$
  block, dictated by the defining relation of $\mathfrak{so}(k,l)$), so
  $\langle X,Y\rangle=\operatorname{tr}(X^*Y)=\operatorname{tr}(XY)$, the negative of the skew case;
  since the Killing form of $\mathfrak{so}(k,l)$ is again $B(X,Y)=(k+l-2)\operatorname{tr}(XY)$
  (Exercise 10.5.6(6)), $\langle X,Y\rangle=\tfrac1{k+l-2}B(X,Y)$, with the opposite
  sign relative to \cref{ex:pet-ch10-compact-grassmannian}. Feeding this into
  \cref{thm:pet-ch10-curvature-bracket-formula} negates $\operatorname{Ric}$ and the
  sectional curvature formula relative to the compact case, giving nonpositive
  curvature. Irreducibility of the isotropy representation again forces, via
  \cref{cor:pet-ch10-maximal-subalgebra}, $\mathrm{iso}=\mathfrak{so}(k,l)$.
\end{proof}

\begin{example}[§10.2.3 — complex projective space as a symmetric space]
  \label{ex:pet-ch10-complex-projective-space}
  \lean{PetersenLib.complexProjectiveSpaceSymmetricSpace}
  \uses{def:pet-ch10-tp-decomposition,thm:pet-ch10-curvature-bracket-formula}
  Let $M=\mathbb{CP}^n=G_1(\mathbb C^{n+1})$, the complex lines in $\mathbb C^{n+1}$, with
  $p=\mathbb C$ the first coordinate axis. Then $\mathrm{SU}(n+1)$ acts transitively on
  $M$ with isotropy $\mathrm S(\mathrm U(1)\times\mathrm U(n))\cong\mathrm U(n)$ via
  $A\mapsto\begin{pmatrix}\det A^{-1}&0\\0&A\end{pmatrix}$, and $M$ is symmetric
  with involution $\begin{pmatrix}(-1)^n&0\\0&-I_n\end{pmatrix}$. At the Lie
  algebra level, $\mathfrak{su}(n+1)=\mathfrak t_p\oplus\mathfrak u(n)$ with
  \[
    \mathfrak t_p=\Big\{\begin{pmatrix}0&-z^*\\z&0\end{pmatrix}\mid z\in\mathbb C^n\Big\},
  \]
  $\langle X,Y\rangle=\operatorname{tr}(X^*Y)=-\operatorname{tr}(XY)=-\tfrac1{2(n+1)}B(X,Y)$ on
  $\mathfrak t_p$, and
  \[
    \operatorname{Ric}=(n+1)g,\qquad
    \langle R(X,Y)Y,X\rangle=|[X,Y]|^2=4|\operatorname{Im}(z^*w)|^2-2\operatorname{Re}(z^*w)^2+2|z|^2|w|^2.
  \]
  On an orthonormal $2$-plane ($|z|^2=|w|^2=\tfrac12$, $\operatorname{Re}(z^*w)=0$) this
  equals $6|z^*w|^2+\tfrac12$, so $\tfrac14\le\sec\le2$: the Fubini–Study metric on
  $\mathbb{CP}^n$ is quarter-pinched, with equality $\sec=2$ iff $w=iz$ and
  $\sec=\tfrac14$ iff $z^*w=0$.
  \source{petersen:page-0393,petersen:page-0394,petersen:page-0395}
\end{example}
\begin{proof}
  \uses{ex:pet-ch10-complex-projective-space}
  $\mathrm U(n+1)$ preserves the Hermitian form $z^*w=\sum\bar z_iw_i$ and acts on
  $M$; scalars $aI$ with $a\bar a=1$ act trivially, so the action factors through
  $\mathrm{SU}(n+1)$, still transitive (with finite kernel $\{aI:a^{n+1}=1\}$). The
  isotropy of $p=\mathbb C$ is the block-diagonal $\mathrm S(\mathrm U(1)\times\mathrm U(n))$,
  identified with $\mathrm U(n)$ as stated; the displayed matrix, having
  determinant $(-1)^n\cdot(-1)^n=1$ and squaring to $I$, fixes $p$ and acts as $-1$
  on the complementary directions, giving the symmetric-space involution.

  Passing to Lie algebras, $\mathfrak{su}(n+1)=\{A=-A^*,\operatorname{tr}A=0\}$,
  $\mathfrak u(n)=\{B=-B^*\}$ includes via $B\mapsto\begin{pmatrix}-\operatorname{tr}B&0\\0&B\end{pmatrix}$,
  so a general element of $\mathfrak{su}(n+1)$ is
  $\begin{pmatrix}-\operatorname{tr}B&-z^*\\z&B\end{pmatrix}$, giving the stated
  $\mathfrak t_p$. For $X,Y\in\mathfrak t_p$, $\langle X,Y\rangle=\operatorname{tr}(X^*Y)=-\operatorname{tr}(XY)$
  directly from the block form; the Killing form of $\mathfrak{su}(m)$ is
  $B(X,Y)=2m\operatorname{tr}(XY)$ (Exercise 10.5.6(3)–(4) with $m=n+1$), giving
  $\langle X,Y\rangle=-\tfrac1{2(n+1)}B(X,Y)$.

  By \cref{thm:pet-ch10-curvature-bracket-formula}, $\operatorname{Ric}=(n+1)g$ and
  $\langle R(X,Y)Y,X\rangle=|[X,Y]|^2$; a direct block computation of $[X,Y]$ for
  $X,Y\in\mathfrak t_p$ (parametrized by $z,w\in\mathbb C^n$) gives
  $[X,Y]=\begin{pmatrix}-(z^*w-w^*z)&0\\0&wz^*-zw^*\end{pmatrix}$, whose squared
  norm expands (using $z^*w-w^*z=2i\operatorname{Im}(z^*w)$ and
  $\operatorname{tr}((wz^*-zw^*)^2)=2|z|^2|w|^2-2|z^*w|^2$) to
  $4|\operatorname{Im}(z^*w)|^2-2\operatorname{Re}(z^*w)^2+2|z|^2|w|^2$. Restricting to an
  orthonormal $2$-plane forces $|z|^2=|w|^2=\tfrac12$, $\operatorname{Re}(z^*w)=0$, so
  $z^*w=i\operatorname{Im}(z^*w)$ and the expression simplifies to
  $6|z^*w|^2+\tfrac12\in[\tfrac14,2]$, attaining $2$ at $w=iz$ and $\tfrac14$ at
  $z^*w=0$.
\end{proof}

\begin{example}[§10.2.4 — $\mathrm{SL}(n)/\mathrm{SO}(n)$]
  \label{ex:pet-ch10-sl-so-quotient}
  \lean{PetersenLib.slSoSymmetricSpace}
  \uses{def:pet-ch10-tp-decomposition,thm:pet-ch10-curvature-bracket-formula}
  Let $\mathfrak{sl}(n)=\{X\in\operatorname{Mat}_{n\times n}\mid\operatorname{tr}X=0\}$, split as
  $\mathfrak{sl}(n)=\mathfrak t\oplus\mathfrak{so}(n)$ into symmetric and skew-symmetric
  traceless matrices, with involution $\sigma(X)=-X^\top$ ($-1$ on $\mathfrak t$,
  $+1$ on $\mathfrak{so}(n)$). On $\mathfrak t$, with the Euclidean metric
  $\langle X,Y\rangle=\operatorname{tr}(X^*Y)=\operatorname{tr}(XY)=\tfrac1{2n}B(X,Y)$, one has
  $\operatorname{Ric}=-ng$ and $\langle R(X,Y)Z,W\rangle=\langle[X,Y],[Z,W]\rangle$; in
  particular $\mathrm{SL}(n)/\mathrm{SO}(n)$ has nonpositive sectional curvature.
  \source{petersen:page-0395,petersen:page-0396}
\end{example}
\begin{proof}
  \uses{ex:pet-ch10-sl-so-quotient}
  $\mathfrak{sl}(n)$ splits as symmetric $\oplus$ skew-symmetric traceless
  matrices since every traceless matrix is the sum of its symmetric and
  skew-symmetric parts, each traceless (a symmetric traceless part has trace
  $0$ automatically, and skew matrices are automatically traceless). The map
  $\sigma(X)=-X^\top$ is a Lie algebra involution ($\sigma[X,Y]=-[X,Y]^\top=[-X^\top,-Y^\top]=[\sigma X,\sigma Y]$
  by the antisymmetry of transpose on the commutator) fixing $\mathfrak{so}(n)$ and
  negating $\mathfrak t$. On $\mathfrak t$, $\langle X,Y\rangle=\operatorname{tr}(X^*Y)=\operatorname{tr}(XY)$
  since $X,Y$ are symmetric real, and the Killing form of $\mathfrak{sl}(n)$
  restricted to $\mathfrak t$ is $2n\operatorname{tr}(XY)$ (Exercise 10.5.6(3)), giving
  $\langle X,Y\rangle=\tfrac1{2n}B(X,Y)$. By
  \cref{thm:pet-ch10-curvature-bracket-formula}, $\operatorname{Ric}=-ng$ (negative sign
  because $\langle X,Y\rangle$ is the \emph{positive} multiple $\tfrac1{2n}B$ rather
  than a negative one, reversing the sign relative to the compact
  \cref{ex:pet-ch10-compact-grassmannian} computation), and
  $\langle R(X,Y)Z,W\rangle=\langle[X,Y],[Z,W]\rangle$ is a sum of squares up to
  sign, giving nonpositive sectional curvature $\langle R(X,Y)Y,X\rangle=-\langle[X,Y],[X,Y]\rangle\le0$.
\end{proof}

\begin{example}[§10.2.5 — Lie groups as symmetric spaces]
  \label{ex:pet-ch10-lie-group-symmetric}
  \lean{PetersenLib.compactLieGroupSymmetricSpace,PetersenLib.noncompactDualSymmetricSpace}
  \uses{rem:pet-ch10-symmetric-space-from-involution,def:pet-ch10-killing-form,thm:pet-ch10-curvature-bracket-formula}
  Let $G$ be a compact Lie group with biinvariant metric and centerless Lie
  algebra $\mathfrak g\cong T_eG$; since left-invariant fields are Killing fields,
  $\mathrm{ad}_X=2(\nabla X)|_e$ is skew, so the Killing form $B$ is negative
  semidefinite, vanishing only on the center, and $g=-B$ recovers the biinvariant
  metric. The Lie algebra description of $G$ as a symmetric space is
  $(\mathfrak g\oplus\mathfrak g,\sigma)$, $\sigma(X,Y)=(Y,X)$: the diagonal
  $\mathfrak g^\Delta$ is the $1$-eigenspace $\mathfrak k$, the antidiagonal
  $\mathfrak g^+=\{(X,-X)\}$ is the $(-1)$-eigenspace $\mathfrak t\cong\mathfrak g$,
  and $G=(G\times G)/G^\Delta$. On $\mathfrak t$ (scaled by a factor of $2$ from
  the biinvariant metric on $\mathfrak g$), $\operatorname{Ric}=-\tfrac12B=\tfrac12g$ and the
  curvature operator is nonnegative. Dually, for the same centerless compact
  $\mathfrak g$, the complexification $(\mathfrak g\otimes\mathbb C,\sigma)$ with
  $\sigma(X)=\bar X$ has $\mathfrak k=\mathfrak g$, $\mathfrak t=i\mathfrak g$, metric
  $g(iX,iY)=-B(X,Y)=B(iX,iY)$ on $\mathfrak t$, and gives a noncompact symmetric
  space with $\operatorname{Ric}(iX,iY)=-\tfrac12g(iX,iY)$, i.e.\ nonpositive curvature.
  \source{petersen:page-0389,petersen:page-0396}
\end{example}
\begin{proof}
  \uses{ex:pet-ch10-lie-group-symmetric}
  \emph{Compact case.} That $\mathrm{ad}_X$ is skew for left-invariant $X$ follows
  because such $X$ are Killing for the biinvariant metric, so $\nabla X$ is
  skew-symmetric, and $\mathrm{ad}_X=2(\nabla X)|_e$ on a biinvariant Lie group.
  Then $B(X,X)=\operatorname{tr}(\mathrm{ad}_X^2)\le0$ as in
  \cref{rem:pet-ch10-adjoint-skew-symmetric}, with equality exactly on the center;
  when $\mathfrak g$ has trivial center, $B$ is negative definite and $g=-B$ is a
  biinvariant metric (since $\mathrm{ad}_X$ is $B$-skew for every Lie algebra).
  Applying \cref{rem:pet-ch10-symmetric-space-from-involution} to
  $\mathfrak g\oplus\mathfrak g$ with $\sigma(X,Y)=(Y,X)$: this is visibly a Lie
  algebra involution, its $1$-eigenspace is the diagonal $\mathfrak g^\Delta\cong\mathfrak g$
  (matching the compact group $K=G^\Delta\cong G$) and its $(-1)$-eigenspace is
  $\mathfrak g^+\cong\mathfrak g$; taking $G'=G\times G$ realizes the recipe of
  \cref{rem:pet-ch10-symmetric-space-from-involution}, giving $G=(G\times G)/G^\Delta$
  as a symmetric space. The Ricci and curvature-sign formulas follow from
  \cref{thm:pet-ch10-curvature-bracket-formula} applied to this description,
  using $B_{\mathfrak g\oplus\mathfrak g}((X,X),(Y,Y))=2B_{\mathfrak g}(X,Y)$ restricted to
  the diagonal-complement identification of $\mathfrak t$ with $\mathfrak g$.

  \emph{Noncompact dual.} On $\mathfrak g\otimes\mathbb C=\mathfrak g\oplus i\mathfrak g$,
  complex conjugation $\sigma(X)=\bar X$ is a Lie algebra involution (the bracket
  on $\mathfrak g\otimes\mathbb C$ is $\mathbb C$-bilinear, hence commutes with the
  real-linear conjugation) fixing $\mathfrak g$ and negating $i\mathfrak g$. The
  Killing form of $\mathfrak g\otimes\mathbb C$ restricted to $\mathfrak g$ is (the
  complexification of) $B$, negative definite by hypothesis, so restricted to
  $i\mathfrak g$ it is positive definite: $g(iX,iY):=-B(X,Y)=B(iX,iY)$ defines a
  genuine (positive-definite) metric on $\mathfrak t=i\mathfrak g$. Applying
  \cref{thm:pet-ch10-curvature-bracket-formula} to this description gives
  $\operatorname{Ric}(iX,iY)=-\tfrac12B(iX,iY)=-\tfrac12g(iX,iY)$, i.e.\ nonpositive
  curvature.
\end{proof}

\section{Holonomy}

\subsection{The Holonomy Group}

\begin{definition}[holonomy group and restricted holonomy group]
  \label{def:pet-ch10-holonomy-group}
  \lean{PetersenLib.HolonomyGroup,PetersenLib.RestrictedHolonomyGroup}
  Let $(M,g)$ be a Riemannian $n$-manifold, $p\in M$. For a unit-speed
  (piecewise smooth) curve $c:[a,b]\to M$ let
  $P_{c(a)}^{c(b)}:T_{c(a)}M\to T_{c(b)}M$ denote parallel translation along $c$.
  For a loop $c(a)=c(b)=p$, $P_p^p$ is an isometry of $T_pM$; the set of all such
  isometries, over all piecewise smooth loops at $p$, forms the \emph{holonomy
  group} $\mathrm{Hol}_p=\mathrm{Hol}_p(M,g)\subset O(T_pM)\cong O(n)$, a (possibly
  noncompact) Lie group. The \emph{restricted holonomy group}
  $\mathrm{Hol}_p^0=\mathrm{Hol}_p^0(M,g)$ is the connected normal subgroup obtained
  using only contractible loops; it is compact and a Lie group.
  \source{petersen:page-0396,petersen:page-0397}
\end{definition}

\begin{remark}[elementary properties of the holonomy group]
  \label{rem:pet-ch10-holonomy-properties}
  \lean{PetersenLib.holonomyEuclidean,PetersenLib.holonomySphere,PetersenLib.holonomyHyperbolicSpace,PetersenLib.holonomyOrientability,PetersenLib.holonomyUniversalCover,PetersenLib.holonomyProduct,PetersenLib.holonomyConjugate,PetersenLib.holonomyInvariantTensor}
  \uses{def:pet-ch10-holonomy-group}
  \begin{enumerate}
    \item[(a)] $\mathrm{Hol}_p(\mathbb R^n)=\{I\}$.
    \item[(b)] $\mathrm{Hol}_p(S^n(R))=\mathrm{SO}(n)$.
    \item[(c)] $\mathrm{Hol}_p(H^n)=\mathrm{SO}(n)$.
    \item[(d)] $\mathrm{Hol}_p(M,g)\subset\mathrm{SO}(n)$ iff $M$ is orientable.
    \item[(e)] $\mathrm{Hol}_p(\widetilde M,\widetilde g)=\mathrm{Hol}_p^0(M,g)$ for the
      universal cover $\widetilde M$.
    \item[(f)] $\mathrm{Hol}_{(p,q)}(M_1\times M_2,g_1+g_2)=\mathrm{Hol}_p(M_1,g_1)\times\mathrm{Hol}_q(M_2,g_2)$.
    \item[(g)] $\mathrm{Hol}_p(M,g)$ is conjugate to $\mathrm{Hol}_q(M,g)$ via parallel
      translation along any curve from $p$ to $q$.
    \item[(h)] A tensor at $p$ extends to a parallel tensor on $(M,g)$ iff it is
      $\mathrm{Hol}_p(M,g)$-invariant; e.g.\ a $2$-form $\omega$ extends iff
      $\omega(Pv,Pw)=\omega(v,w)$ for all $P\in\mathrm{Hol}_p(M,g)$,
      $v,w\in T_pM$.
  \end{enumerate}
  \source{petersen:page-0397}
\end{remark}

By property (f), a Cartesian-product splitting is reflected in a product structure
of the holonomy, and by (g), if the (restricted) holonomy representation on
$T_pM$ decomposes at one point it decomposes at every point (via parallel
translation), yielding a global decomposition $TM=\eta_1\oplus\cdots\oplus\eta_k$
of the tangent bundle into distributions invariant under parallel translation,
obtained from an orthogonal decomposition of $T_pM$ into irreducible invariant
subspaces for $\mathrm{Hol}_p^0$.

\begin{theorem}[Thm. 10.3.1 — de Rham decomposition (1952)]
  \label{thm:pet-ch10-de-rham-decomposition}
  \lean{PetersenLib.deRhamDecomposition}
  \uses{rem:pet-ch10-holonomy-properties}
  Let $(M,g)$ be a Riemannian manifold and $TM=\eta_1\oplus\cdots\oplus\eta_k$
  the decomposition of the tangent bundle into irreducible components for the
  restricted holonomy. Around each $p\in M$ there is a neighborhood $U$ with a
  product structure $(U,g)=(U_1\times\cdots\times U_k,g_1+\cdots+g_k)$,
  $TU_i=\eta_i|_U$. If $(M,g)$ is simply connected and complete, this splitting
  is global: $(M,g)=(M_1\times\cdots\times M_k,g_1+\cdots+g_k)$, $TM_i=\eta_i$.
  \source{petersen:page-0397,petersen:page-0398}
\end{theorem}
\begin{proof}
  \uses{thm:pet-ch10-de-rham-decomposition}
  Each $\eta_i$ is integrable: being invariant under parallel translation, and
  hence under the Levi-Civita connection, the distribution $\eta_i$ has
  vanishing torsion of the induced connection, so its Lie bracket of sections
  stays in $\eta_i$ (Frobenius integrability). This gives, at the manifold
  level, a local splitting into maximal integral submanifolds through $p$. Each
  integral submanifold is totally geodesic, since its tangent space is invariant
  under parallel translation along curves within it; hence the ambient metric
  restricted to a product neighborhood of integral submanifolds splits as
  $g_1+\cdots+g_k$, giving the local statement.

  For the global statement, let $M_i$ be the maximal integral submanifold of
  $\eta_i$ through a fixed $p\in M$, and consider the abstract product Riemannian
  manifold $(M_1\times\cdots\times M_k,\,g_1+\cdots+g_k)$; by the local splitting,
  $(M,g)$ and this product are isometric near $p$. Since $(M,g)$ is complete and
  each $M_i$ is totally geodesic, $M_i$ (and hence the product) is complete. The
  local isometry extends, by analytic continuation of the local isometric
  embedding using completeness of both sides, to a global isometric immersion
  $(M,g)\to(M_1\times\cdots\times M_k,g_1+\cdots+g_k)$; completeness makes this
  immersion onto and in fact a Riemannian covering map. As $M$ is simply
  connected, $M$ is therefore isometric to the universal cover of the product,
  i.e.\ to $(\widetilde M_1\times\cdots\times\widetilde M_k,\widetilde g_1+\cdots+\widetilde g_k)$
  with the pulled-back product metric.
\end{proof}

\subsection{Rough Classification of Symmetric Spaces}

Guided by the de Rham decomposition, classification questions reduce to the case
of an \emph{irreducible} Riemannian manifold, i.e.\ one whose holonomy has no
invariant subspaces.

\begin{theorem}[Thm. 10.3.2 — parallel tensors on irreducible manifolds]
  \label{thm:pet-ch10-parallel-tensor-eigenvalue}
  \lean{PetersenLib.parallelTensorSingleEigenvalue}
  \uses{thm:pet-ch10-de-rham-decomposition}
  Let $(M,g)$ be an irreducible Riemannian manifold with a parallel $(1,1)$-tensor
  $T$. Then the symmetric part $S=\tfrac12(T+T^*)$ and skew-symmetric part
  $A=\tfrac12(T-T^*)$ each have exactly one (complex) eigenvalue. If $A\ne0$, it
  induces (after scaling) a parallel complex structure, i.e.\ a Kähler structure.
  \source{petersen:page-0399}
\end{theorem}
\begin{proof}
  \uses{thm:pet-ch10-parallel-tensor-eigenvalue}
  Since $g$ is parallel and $T$ is parallel, so is its adjoint: for vector fields
  $X,Y,Z$,
  \[
    g((\nabla_XT)(Y),Z)
    =\nabla_Xg(T(Y),Z)-g(T(Y),\nabla_XZ)-g(T(\nabla_XY),Z)
    =g(Y,(\nabla_XT^*)(Z)),
  \]
  using $\nabla g=0$ to rewrite $g(T(Y),\nabla_XZ)=\nabla_Xg(Y,T^*Z)-g(Y,(\nabla_XT^*)(Z))$
  and matching terms; hence $\nabla T=0$ implies $\nabla T^*=0$, so
  $\nabla S=\nabla A=0$, and both $S,A$ are invariant under parallel translation.

  Decompose $T_pM=E_1\oplus\cdots\oplus E_k$ into the orthogonal eigenspaces of
  $S$ for its distinct (real) eigenvalues $\lambda_1<\cdots<\lambda_k$. Since $S$
  is parallel, this decomposition parallel-translates to a global decomposition
  $TM=\eta_1\oplus\cdots\oplus\eta_k$ with $S|_{\eta_i}=\lambda_iI$; each $\eta_i$
  is preserved by the holonomy (being defined by parallel translation), so
  irreducibility of $(M,g)$ forces $k=1$: $S$ has a single eigenvalue.

  Similarly, $A$ has purely imaginary eigenvalues, giving an orthogonal
  decomposition $T_pM=F_1\oplus\cdots\oplus F_l$ into $A$-invariant subspaces on
  which $A$ has complex eigenvalue $i\mu_\mu$, $\mu_1<\cdots<\mu_l$; the same
  irreducibility argument forces $l=1$. If $\mu_1\ne0$, then $\tfrac1{\mu_1}A$ is
  parallel with only eigenvalue $i$, hence squares to $-I$: a parallel complex
  structure, giving a Kähler structure on $(M,g)$.
\end{proof}

\begin{corollary}[Cor. 10.3.3 — irreducible symmetric spaces are Einstein]
  \label{cor:pet-ch10-irreducible-symmetric-einstein}
  \lean{PetersenLib.irreducibleSymmetricSpaceEinstein}
  \uses{thm:pet-ch10-parallel-tensor-eigenvalue}
  A simply connected irreducible symmetric space is an Einstein manifold; in
  particular its curvature operator is nonnegative or nonpositive according to
  the sign of the Einstein constant.
  \source{petersen:page-0399}
\end{corollary}

\begin{remark}[compact, flat, and noncompact type]
  \label{rem:pet-ch10-compact-flat-noncompact-type}
  \lean{PetersenLib.symmetricSpaceCompactFlatNoncompactType}
  \uses{cor:pet-ch10-irreducible-symmetric-einstein}
  For an irreducible symmetric space, \cref{cor:pet-ch10-irreducible-symmetric-einstein}
  gives an Einstein constant, and three cases occur. \emph{Compact type}: if the
  Einstein constant is positive, Myers' diameter bound forces $M$ compact, with
  nonnegative curvature operator. \emph{Flat type}: if $M$ is Ricci flat, it is
  flat, and the only irreducible Ricci-flat examples are $S^1$ and $\mathbb R^1$.
  \emph{Noncompact type}: if the Einstein constant is negative, Bochner's theorem
  on Killing fields forces $M$ noncompact, with nonpositive curvature operator.
  \source{petersen:page-0399,petersen:page-0400}
\end{remark}

\begin{definition}[Types I--IV of irreducible symmetric spaces]
  \label{def:pet-ch10-symmetric-space-type-I-IV}
  \lean{PetersenLib.symmetricSpaceTypeI,PetersenLib.symmetricSpaceTypeII,PetersenLib.symmetricSpaceTypeIII,PetersenLib.symmetricSpaceTypeIV}
  \uses{rem:pet-ch10-compact-flat-noncompact-type,rem:pet-ch10-symmetric-space-from-involution}
  Irreducible symmetric spaces of compact and noncompact type occur in dual
  pairs described by Lie algebra pairs $(\mathfrak g,\mathfrak t_p)$ with an
  involution having $1$-eigenspace $\mathfrak k$ (the Lie algebra of a compact
  group acting on the $(-1)$-eigenspace $\mathfrak t_p$):
  \begin{itemize}
    \item \textbf{Type I}: compact, $(\mathfrak g,\mathfrak t)$ with $\mathfrak g$
      simple, the Lie algebra of a compact group, $\mathfrak t\subset\mathfrak g$
      a maximal subalgebra; e.g.\ $(\mathfrak{so}(k+l),\mathfrak{so}(k)\times\mathfrak{so}(l))$.
    \item \textbf{Type II}: compact, $(\mathfrak t\oplus\mathfrak t,\Delta\mathfrak t)$
      with $\mathfrak t$ simple, the Lie algebra of a compact Lie group $G$; the
      space is $G$ with a biinvariant metric.
    \item \textbf{Type III}: noncompact, $(\mathfrak g,\mathfrak t)$ with
      $\mathfrak g$ simple, the Lie algebra of a noncompact group,
      $\mathfrak t\subset\mathfrak g$ maximal corresponding to a compact group;
      e.g.\ $(\mathfrak{so}(k,l),\mathfrak{so}(k)\times\mathfrak{so}(l))$ or
      $(\mathfrak{sl}(n),\mathfrak{so}(n))$.
    \item \textbf{Type IV}: noncompact, $(\mathfrak t\otimes\mathbb C,\mathfrak t)$
      with $\mathfrak t$ simple, corresponding to a compact group, and
      $\mathfrak t\otimes\mathbb C=\mathfrak t\oplus i\mathfrak t$ its
      complexification; e.g.\ $(\mathfrak{so}(n,\mathbb C),\mathfrak{so}(n))$.
  \end{itemize}
  In every case \cref{cor:pet-ch10-maximal-subalgebra} shows $t_p=\mathrm{iso}_p$.
  Since compact-type symmetric spaces have nonnegative curvature operator, the
  Bochner technique shows all their harmonic forms are parallel, hence
  holonomy-invariant, reducing the computation of their cohomology to a
  classical invariance problem for the isotropy representation — used, e.g., to
  compute Pontryagin and Chern classes from the cohomology of real and complex
  Grassmannians.
  \source{petersen:page-0400}
\end{definition}

\subsection{Curvature and Holonomy}

\begin{definition}[Berger's list of irreducible holonomy groups]
  \label{def:pet-ch10-berger-holonomy-table}
  \lean{PetersenLib.bergerHolonomyList}
  For a simply connected irreducible Riemannian $n$-manifold whose holonomy
  acts transitively on the unit sphere of $T_pM$, the possible holonomy groups
  $\mathrm{Hol}_p$, and the resulting geometric structure, are:
  \begin{center}
    \begin{tabular}{c|c|c}
      $\dim=n$ & $\mathrm{Hol}_p$ & properties \\\hline
      $n$ & $\mathrm{SO}(n)$ & generic case \\
      $n=2m$ & $\mathrm U(m)$ & Kähler \\
      $n=2m$ & $\mathrm{SU}(m)$ & Kähler and Ricci flat \\
      $n=4m$ & $\mathrm{Sp}(1)\cdot\mathrm{Sp}(m)$ & quaternionic-Kähler and Einstein \\
      $n=4m$ & $\mathrm{Sp}(m)$ & hyper-Kähler and Ricci flat \\
      $n=16$ & $\mathrm{Spin}(9)$ & symmetric and Einstein \\
      $n=8$ & $\mathrm{Spin}(7)$ & Ricci flat \\
      $n=7$ & $G_2$ & Ricci flat
    \end{tabular}
  \end{center}
  \source{petersen:page-0401}
\end{definition}

\begin{theorem}[Thm. 10.3.4 — Berger (1955) and Simons (1962)]
  \label{thm:pet-ch10-berger-simons-classification}
  \lean{PetersenLib.bergerSimonsClassification}
  \uses{def:pet-ch10-berger-holonomy-table,def:pet-ch10-rank}
  Let $(M,g)$ be a simply connected irreducible Riemannian $n$-manifold. Then
  $\mathrm{Hol}_p$ either acts transitively on the unit sphere of $T_pM$, in
  which case $\mathrm{Hol}_p$ is one of the groups of
  \cref{def:pet-ch10-berger-holonomy-table}, or $(M,g)$ is a symmetric space of
  rank $\ge2$.
  \source{petersen:page-0401}
\end{theorem}

\begin{corollary}[Cor. 10.3.5 — nontransitive holonomy forces reducibility or rank $\ge2$]
  \label{cor:pet-ch10-nontransitive-holonomy}
  \lean{PetersenLib.nontransitiveHolonomyReducibleOrSymmetric}
  \uses{thm:pet-ch10-berger-simons-classification}
  If a Riemannian manifold's holonomy does not act transitively on the unit
  sphere, then it is either reducible or a locally symmetric space of rank
  $\ge2$; in particular its rank is $\ge2$.
  \source{petersen:page-0402}
\end{corollary}

\begin{theorem}[Thm. 10.3.6 — rank rigidity for nonpositive curvature (Ballmann; Burns–Spatzier)]
  \label{thm:pet-ch10-rank-rigidity-nonpositive}
  \lean{PetersenLib.rankRigidityNonpositiveCurvature}
  \uses{cor:pet-ch10-nontransitive-holonomy}
  A compact Riemannian manifold of nonpositive curvature and rank $\ge2$ does
  not have transitive holonomy; in particular it is either reducible or locally
  symmetric.
  \source{petersen:page-0402}
\end{theorem}

\begin{theorem}[Thm. 10.3.7 — Gallot and D. Meyer (1975)]
  \label{thm:pet-ch10-gallot-meyer}
  \lean{PetersenLib.gallotMeyerNonnegativeCurvatureOperator}
  \uses{thm:pet-ch10-berger-simons-classification,thm:pet-ch10-parallel-tensor-eigenvalue,thm:pet-ch10-de-rham-decomposition}
  Let $(M,g)$ be a compact Riemannian $n$-manifold with nonnegative curvature
  operator. Then one of:
  \begin{enumerate}
    \item[(a)] $(M,g)$ is reducible or locally symmetric;
    \item[(b)] $\mathrm{Hol}^0(M,g)=\mathrm{SO}(n)$ and the universal cover is a
      homology sphere;
    \item[(c)] $\mathrm{Hol}^0(M,g)=\mathrm U(n/2)$ and the universal cover is a
      homology $\mathbb{CP}^{n/2}$.
  \end{enumerate}
  \source{petersen:page-0402,petersen:page-0403}
\end{theorem}
\begin{proof}
  \uses{thm:pet-ch10-gallot-meyer}
  By the structure theory for nonnegative Ricci curvature, the universal cover
  is isometric to $N\times\mathbb R^k$ with $k>0$ if $\pi_1(M)$ is infinite, so $M$
  is then reducible; hence we may assume $M$ is simply connected. By the Hodge
  theorem, either $H^p(M,\mathbb R)=0$ for all $0<p<n$ (a homology sphere), or some
  $H^p(M,\mathbb R)\ne0$ for $0<p<n$ carries a nonzero harmonic $p$-form $\omega$; as
  the curvature operator is nonnegative, the Bochner technique shows $\omega$ is
  parallel.

  Assume $M$ is irreducible with transitive holonomy (the reducible case is (a),
  and by \cref{thm:pet-ch10-berger-simons-classification} the only remaining
  alternative to transitivity is that $M$ is symmetric, again giving (a)). The
  Ricci-flat holonomies in
  \cref{def:pet-ch10-berger-holonomy-table} are impossible, since nonnegative
  curvature together with Ricci flatness forces $M$ flat. This leaves
  $\mathrm{SO}(n)$, $\mathrm U(n/2)$, $\mathrm{Sp}(1)\mathrm{Sp}(n/4)$, or
  $\mathrm{Spin}(9)$; in the latter two, holonomy considerations force $M$
  Einstein, and Tachibana's theorem (nonnegative curvature operator plus
  Einstein implies locally symmetric) gives (a), with the classification of
  symmetric spaces identifying $M$ with $\mathbb{HP}^{n/4}$ or $\mathbb{OP}^2$.

  Suppose $\mathrm{Hol}_p=\mathrm{SO}(n)$ and $\omega$ is a parallel $p$-form,
  $0<p<n$. For $v_1,\dots,v_p\in T_pM$ there is $P\in\mathrm{SO}(n)$ with
  $P(v_i)=v_i$ for $i\ge2$ and $P(v_1)=-v_1$ (reflect $v_1$, correct the
  determinant using an unused dimension since $p<n$); holonomy invariance of
  $\omega$ then gives $\omega(v_1,\dots,v_p)=\omega(Pv_1,\dots,Pv_p)=-\omega(v_1,\dots,v_p)$,
  so $\omega=0$, and case (b) holds by exclusion (there is no nonzero
  intermediate harmonic form, so $M$ is a homology sphere).

  Suppose instead $\mathrm{Hol}_p=\mathrm U(n/2)$, i.e.\ $(M,g)$ is Kähler. The
  invariant parallel complex structure gives, after type change, a parallel
  $2$-form $\omega$; by \cref{thm:pet-ch10-parallel-tensor-eigenvalue}, any other
  parallel $2$-form is a multiple of $\omega$, so $\dim H^2=1$. Odd cohomology
  vanishes as in the $\mathrm{SO}(n)$ case, using $P=-I\in\mathrm U(n/2)$ (valid
  since $n$ is even). More generally, for a $\mathrm U(n/2)$-invariant $p$-form
  $\theta$, $1<p<n$, and an orthonormal basis $e_1,\dots,e_n$, one finds
  $P\in\mathrm U(n/2)$ fixing $e_1,\dots,e_{p-1}$ and sending $e_p\mapsto e_{i_p}$
  for any index $i_p$, so $\theta(e_1,\dots,e_p)=\theta(e_1,\dots,e_{p-1},e_{i_p})$;
  iterating shows $\theta$ is constant on ordered $p$-tuples of basis vectors, so
  all $p$-forms are proportional. The powers $\omega^k=\omega\wedge\cdots\wedge\omega$
  are nonzero parallel forms, hence generate $H^{2k}$, giving $M$ the cohomology
  ring of $\mathbb{CP}^{n/2}$: case (c).
\end{proof}

\begin{theorem}[Thm. 10.3.8 — Brendle and Schoen (2008)]
  \label{thm:pet-ch10-brendle-schoen}
  \lean{PetersenLib.brendleSchoenNonnegativeComplexSectionalCurvature}
  \uses{thm:pet-ch10-gallot-meyer}
  Let $(M,g)$ be a compact Riemannian $n$-manifold with nonnegative complex
  sectional curvature. Then one of:
  \begin{enumerate}
    \item[(a)] $(M,g)$ is reducible or locally symmetric;
    \item[(b)] $M$ is diffeomorphic to a space of constant positive curvature;
    \item[(c)] the universal cover of $M$ is biholomorphic to $\mathbb{CP}^{n/2}$.
  \end{enumerate}
  \source{petersen:page-0404}
\end{theorem}

\section{Further Study}

Petersen points the reader to Besse, Helgason, Jost, Eberlein, O'Neill, and
Klingenberg for topics on symmetric spaces beyond the scope of this chapter,
notably the fine classification of the irreducible symmetric spaces of each
type and the theory of Lie groups underlying it.

\section{Exercises}

\begin{remark}[Exercise 10.5.1 — geodesic loops are closed geodesics]
  \label{rem:pet-ch10-exercise-symmetric-flow-axis}
  \lean{PetersenLib.exercise10_5_1}
  \uses{def:pet-ch10-tp-decomposition}
  Let $M$ be a symmetric space and $X\in t_p$ nontrivial (so $X$ is a Killing
  field with $(\nabla X)|_p=0$). Show: (1) the flow of $X$ is
  $F_s=A_{\exp(sX|_p)}\circ A_p$; (2) $c(t)=\exp(tX|_p)$ is an axis for $F_s$,
  i.e.\ $F_s(c(t))=c(t+s)$; (3) if $c(a)=c(b)$ then $\dot c(a)=\dot c(b)$;
  (4) conclude that geodesic loops are always closed geodesics.
  \source{petersen:page-0405}
\end{remark}

\begin{remark}[Exercise 10.5.2 — the fundamental group of a symmetric space is abelian]
  \label{rem:pet-ch10-exercise-pi1-abelian}
  \lean{PetersenLib.exercise10_5_2}
  \uses{rem:pet-ch10-exercise-symmetric-flow-axis}
  Let $M$ be a symmetric space. Show that $A_p$ acting on $\pi_1(M,p)$ by
  $[\gamma]\mapsto[\gamma]^{-1}$ shows $\pi_1(M,p)$ is abelian. (Hint: every
  element of $\pi_1(M,p)$ is represented by a geodesic loop, which by
  \cref{rem:pet-ch10-exercise-symmetric-flow-axis} is a closed geodesic.)
  \source{petersen:page-0405}
\end{remark}

\begin{remark}[Exercise 10.5.3 — Jacobi fields along a geodesic in a symmetric space]
  \label{rem:pet-ch10-exercise-jacobi-fields-eigenvector}
  \lean{PetersenLib.exercise10_5_3}
  \uses{def:pet-ch10-tp-decomposition}
  Let $M$ be a symmetric space and $c$ a geodesic. (1) If $E(t)$ is parallel
  along $c$, show $R(E,\dot c)\dot c$ is parallel. (2) If $E(0)$ is an eigenvector
  of $R(\cdot,\dot c(0))\dot c(0)$ with eigenvalue $\kappa$, show
  $J_1(t)=\operatorname{sn}_\kappa(t)E(t)$ and $J_2(t)=\operatorname{sn}_\kappa'(t)E(t)$ are Jacobi
  fields along $c$. (3) If $J$ is a Jacobi field with $J(0)=0$, $J'(t_0)=0$, show
  $J(t_0/u)=0$ for the appropriate rescaling parameter $u$. (4) With $J$ as in
  (3), construct a geodesic variation $c(s,t)$ with $\partial_sc(0,t)=J(t)$,
  $c(s,0)=c(0)$, $c(s,t_0)=c(t_0)$.
  \source{petersen:page-0405,petersen:page-0406}
\end{remark}

\begin{remark}[Exercise 10.5.4 — nonnegative curvature of compact symmetric spaces]
  \label{rem:pet-ch10-exercise-compact-nonneg-curvature}
  \lean{PetersenLib.exercise10_5_4}
  \uses{rem:pet-ch10-exercise-jacobi-fields-eigenvector}
  Let $M$ be a symmetric space. (1) Show directly that if $M$ is compact then
  $\sec\ge0$ (hint: argue by contradiction, producing an unbounded Jacobi field
  along a geodesic). (2) Show that if $c$ is a closed geodesic then
  $R(\cdot,\dot c)\dot c$ has no negative eigenvalues. (3) Show that if
  $\operatorname{Ric}>0$ then $\sec\ge0$.
  \source{petersen:page-0406}
\end{remark}

\begin{remark}[Exercise 10.5.5 — equivalent characterizations of a flat parallel direction]
  \label{rem:pet-ch10-exercise-jacobi-field-equivalence}
  \lean{PetersenLib.exercise10_5_5}
  \uses{def:pet-ch10-holonomy-group}
  Assume $M$ has nonpositive or nonnegative sectional curvature, $c$ a geodesic,
  $E$ a parallel field along $c$. Show the following are equivalent: (1)
  $g(R(E,\dot c)\dot c,E)=0$ everywhere along $c$; (2) $R(E,\dot c)\dot c=0$
  everywhere along $c$; (3) $E$ is a Jacobi field.
  \source{petersen:page-0405,petersen:page-0406}
\end{remark}

\begin{remark}[Exercise 10.5.6 — Killing forms of the classical Lie algebras]
  \label{rem:pet-ch10-exercise-killing-form-classical}
  \lean{PetersenLib.exercise10_5_6}
  \uses{def:pet-ch10-killing-form}
  (1) Show the Killing form of $\mathfrak g\otimes\mathbb C$ is the
  complexification of the Killing form $B$ of a real Lie algebra $\mathfrak g$.
  (2) Show the Killing forms of $\mathfrak{gl}(n,\mathbb C)$ and $\mathfrak{gl}(n)$
  are $B(X,Y)=2n\operatorname{tr}(XY)-2\operatorname{tr}X\operatorname{tr}Y$ (hint: use the basis
  $E_{ij}=[\delta_{i\delta_j}]$). (3) Show that on
  $\mathfrak{sl}(n,\mathbb C)=\mathfrak{sl}(n)\otimes\mathbb C$, $B(X,Y)=2n\operatorname{tr}(XY)$
  (hint: use (2) and that $I\in\mathfrak{gl}(n,\mathbb C)$ commutes with
  $\mathfrak{sl}(n,\mathbb C)$). (4) Show $\mathfrak{sl}(k+l,\mathbb C)=\mathfrak{su}(k,l)\otimes\mathbb C$,
  and conclude $\mathfrak{sl}(k+l)$ and $\mathfrak{su}(k,l)$ have Killing form
  $B(X,Y)=2(k+l)\operatorname{tr}(XY)$. (5) Show that on
  $\mathfrak{so}(n,\mathbb C)=\mathfrak{so}(n)\otimes\mathbb C$, $B(X,Y)=(n-2)\operatorname{tr}(XY)$
  (hint: use the basis $E_{ij}-E_{ji}$, $i<j$). (6) Show
  $\mathfrak{so}(k+l,\mathbb C)=\mathfrak{so}(k,l)\otimes\mathbb C$, and conclude
  $\mathfrak{so}(k,l)$ has Killing form $B(X,Y)=(k+l-2)\operatorname{tr}(XY)$.
  \source{petersen:page-0406}
\end{remark}

\begin{remark}[Exercise 10.5.7 — nondegenerate bilinear forms as a symmetric space]
  \label{rem:pet-ch10-exercise-gl-so-symmetric}
  \lean{PetersenLib.exercise10_5_7}
  \uses{rem:pet-ch10-symmetric-space-from-involution}
  Show that $\mathrm{GL}^+(p+q,\mathbb R)/\mathrm{SO}(p,q)$ defines a symmetric space,
  identifiable with the nondegenerate bilinear forms on $\mathbb R^{p+q}$ of index
  $q$.
  \source{petersen:page-0406}
\end{remark}

\begin{remark}[Exercise 10.5.8 — $\mathrm U(p,q)/\mathrm{SO}(p,q)$ as a symmetric space]
  \label{rem:pet-ch10-exercise-u-pq-so-pq-symmetric}
  \lean{PetersenLib.exercise10_5_8}
  \uses{rem:pet-ch10-symmetric-space-from-involution}
  Show that $\mathrm U(p,q)/\mathrm{SO}(p,q)$ defines a symmetric space and that
  $\mathfrak u(p,q)\otimes\mathbb C=\mathfrak{gl}(p+q,\mathbb C)$.
  \source{petersen:page-0406}
\end{remark}

\begin{remark}[Exercise 10.5.9 — holonomy of $\mathbb{CP}^n$]
  \label{rem:pet-ch10-exercise-cpn-holonomy}
  \lean{PetersenLib.exercise10_5_9}
  \uses{ex:pet-ch10-complex-projective-space,def:pet-ch10-holonomy-group}
  Show that the holonomy of $\mathbb{CP}^n$ is $\mathrm U(n)$.
  \source{petersen:page-0406}
\end{remark}

\begin{remark}[Exercise 10.5.10 — covering spaces of symmetric spaces]
  \label{rem:pet-ch10-exercise-covering-symmetric}
  \lean{PetersenLib.exercise10_5_10}
  \uses{prop:pet-ch10-sigma-p-involution}
  Show that a covering space of a symmetric space is a symmetric space; give an
  example showing the converse fails.
  \source{petersen:page-0406}
\end{remark}

\begin{remark}[Exercise 10.5.11 — flatness iff discrete holonomy]
  \label{rem:pet-ch10-exercise-flat-discrete-holonomy}
  \lean{PetersenLib.exercise10_5_11}
  \uses{def:pet-ch10-holonomy-group}
  Show that a manifold is flat if and only if the (restricted) holonomy Lie
  algebra $\mathfrak{hol}_p=\{0\}$, i.e.\ the holonomy group is discrete.
  \source{petersen:page-0406}
\end{remark}

\begin{remark}[Exercise 10.5.12 — finite fundamental group under irreducible holonomy and $\operatorname{Ric}\ge0$]
  \label{rem:pet-ch10-exercise-finite-fundamental-group}
  \lean{PetersenLib.exercise10_5_12}
  \uses{def:pet-ch10-holonomy-group}
  Show that a compact Riemannian manifold with irreducible restricted holonomy
  and $\operatorname{Ric}\ge0$ has finite fundamental group.
  \source{petersen:page-0406}
\end{remark}

\begin{remark}[Exercise 10.5.13 — identifying $\mathrm{SL}(2,\mathbb R)/\mathrm{SO}(2)$ and $\mathrm{SL}(2,\mathbb C)/\mathrm{SU}(2)$]
  \label{rem:pet-ch10-exercise-sl2-quotients}
  \lean{PetersenLib.exercise10_5_13}
  \uses{ex:pet-ch10-sl-so-quotient}
  Which known spaces are described by $\mathrm{SL}(2,\mathbb R)/\mathrm{SO}(2)$ and
  $\mathrm{SL}(2,\mathbb C)/\mathrm{SU}(2)$?
  \source{petersen:page-0406}
\end{remark}

\begin{remark}[Exercise 10.5.14 — Kähler structure iff holonomy in $\mathrm U(m)$]
  \label{rem:pet-ch10-exercise-kahler-unitary-holonomy}
  \lean{PetersenLib.exercise10_5_14}
  \uses{def:pet-ch10-holonomy-group,def:pet-ch10-berger-holonomy-table}
  Show that the holonomy of a Riemannian manifold is contained in $\mathrm U(m)$
  if and only if it has a Kähler structure.
  \source{petersen:page-0406}
\end{remark}

\begin{remark}[Exercise 10.5.15 — maximal isotropy forces constant curvature]
  \label{rem:pet-ch10-exercise-isop-so-constant-curvature}
  \lean{PetersenLib.exercise10_5_15}
  \uses{def:pet-ch10-tp-decomposition}
  Show that if a homogeneous space has $\mathrm{iso}_p=\mathfrak{so}(T_pM)$ at
  some $p$, then it has constant curvature.
  \source{petersen:page-0406}
\end{remark}

\begin{remark}[Exercise 10.5.16 — maximality of two subalgebras of $\mathfrak{so}(n)$]
  \label{rem:pet-ch10-exercise-maximal-subalgebras-son}
  \lean{PetersenLib.exercise10_5_16}
  \uses{ex:pet-ch10-compact-grassmannian,ex:pet-ch10-complex-projective-space}
  Show that the subalgebras $\mathfrak{so}(k)\times\mathfrak{so}(n-k)$ and
  $\mathfrak u(n/2)$ are maximal in $\mathfrak{so}(n)$.
  \source{petersen:page-0406}
\end{remark}

\begin{remark}[Exercise 10.5.17 — $\mathfrak{su}(m)\subset\mathfrak{so}(m^2-1)$]
  \label{rem:pet-ch10-exercise-sum-in-so}
  \lean{PetersenLib.exercise10_5_17}
  \uses{def:pet-ch10-killing-form}
  Show that $\mathfrak{su}(m)\subset\mathfrak{so}(m^2-1)$. (Hint: let
  $\mathfrak{su}(m)$ act on itself by the adjoint representation.)
  \source{petersen:page-0406}
\end{remark}

\begin{remark}[Exercise 10.5.18 — $t_p\subset\mathfrak{hol}_p$ in general]
  \label{rem:pet-ch10-exercise-tp-subset-holp}
  \lean{PetersenLib.exercise10_5_18}
  \uses{def:pet-ch10-tp-decomposition,def:pet-ch10-holonomy-group}
  Show that for any Riemannian manifold $t_p\subset\mathfrak{hol}_p$ (the
  holonomy Lie algebra). Give an example where equality fails.
  \source{petersen:page-0407}
\end{remark}

\begin{remark}[Exercise 10.5.19 — $t_p=\mathfrak{hol}_p$ for symmetric spaces]
  \label{rem:pet-ch10-exercise-tp-equals-holp-symmetric}
  \lean{PetersenLib.exercise10_5_19}
  \uses{rem:pet-ch10-exercise-tp-subset-holp,cor:pet-ch10-maximal-subalgebra}
  Show that for a symmetric space, $t_p=\mathfrak{hol}_p$. Use this to show
  that, unless the curvature is constant, $t_p=\mathfrak{hol}_p=\mathrm{iso}_p$
  provided $t_p\subset\mathfrak{so}(T_pM)$ is maximal.
  \source{petersen:page-0407}
\end{remark}

\begin{remark}[Exercise 10.5.20 — nonpositive curvature operator of two further quotients]
  \label{rem:pet-ch10-exercise-son-c-sln-c-symmetric}
  \lean{PetersenLib.exercise10_5_20}
  \uses{ex:pet-ch10-lie-group-symmetric}
  Show that $\mathrm{SO}(n,\mathbb C)/\mathrm{SO}(n)$ and
  $\mathrm{SL}(n,\mathbb C)/\mathrm{SU}(n)$ are symmetric spaces with nonpositive
  curvature operator.
  \source{petersen:page-0407}
\end{remark}

\begin{remark}[Exercise 10.5.21 — quaternionic projective space as $\mathrm{Sp}(n+1)/(\mathrm{Sp}(1)\times\mathrm{Sp}(n))$]
  \label{rem:pet-ch10-exercise-quaternionic-projective-space}
  \lean{PetersenLib.exercise10_5_21}
  \uses{ex:pet-ch10-complex-projective-space}
  The \emph{quaternionic projective space} $\mathbb{HP}^n$ is the space of
  quaternionic lines in $\mathbb H^{n+1}$. Define $\mathrm{Sp}(n)\subset\mathrm{SU}(2n)\subset\mathrm{SO}(4n)$
  as the orthogonal matrices commuting with the three complex structures
  generated by $i,j,k$ on $\mathbb R^{4n}$, equivalently the $n\times n$
  quaternionic matrices $A$ with $A^{-1}=A^*$. Show that, viewing $\mathbb H^{n+1}$
  as a right (or left) $\mathbb H$-module, $\mathbb{HP}^n=\mathrm{Sp}(n+1)/(\mathrm{Sp}(1)\times\mathrm{Sp}(n))$.
  \source{petersen:page-0407}
\end{remark}

\begin{remark}[Exercise 10.5.22 — hyperbolic analogues of complex projective space]
  \label{rem:pet-ch10-exercise-hyperbolic-complex-projective}
  \lean{PetersenLib.exercise10_5_22}
  \uses{ex:pet-ch10-complex-projective-space,ex:pet-ch10-hyperbolic-grassmannian}
  Construct the hyperbolic analogues of the complex projective spaces. Show
  that they have negative curvature and are quarter-pinched.
  \source{petersen:page-0407}
\end{remark}

\begin{remark}[Exercise 10.5.23 — the monodromy map of a locally symmetric space]
  \label{rem:pet-ch10-exercise-lie-algebra-locally-symmetric}
  \lean{PetersenLib.exercise10_5_23}
  \uses{rem:pet-ch10-symmetric-space-from-involution}
  Give a Lie algebra description of a locally symmetric space (not necessarily
  complete), and explain why it corresponds to a global symmetric space.
  Conclude that a simply connected locally symmetric space admits a monodromy
  map into a unique global symmetric space, bijective when the locally
  symmetric space is complete.
  \source{petersen:page-0407}
\end{remark}

\begin{remark}[Exercise 10.5.24 — strict curvature-operator sign forces constant curvature]
  \label{rem:pet-ch10-exercise-pinched-curvature-operator-constant}
  \lean{PetersenLib.exercise10_5_24}
  \uses{cor:pet-ch10-irreducible-symmetric-einstein}
  Show that if an irreducible symmetric space has strictly positive or strictly
  negative curvature operator, then it has constant curvature.
  \source{petersen:page-0407}
\end{remark}

\begin{remark}[Exercise 10.5.25 — $[t_p,\mathrm{iso}_p]\subset t_p$]
  \label{rem:pet-ch10-exercise-bracket-tp-isop}
  \lean{PetersenLib.exercise10_5_25}
  \uses{prop:pet-ch10-lie-bracket-curvature}
  Let $M$ be a symmetric space. Show that if $X\in t_p$ and $Y\in\mathrm{iso}_p$,
  then $[X,Y]\in t_p$.
  \source{petersen:page-0407}
\end{remark}

\begin{remark}[Exercise 10.5.26 — curvature and Ricci tensor via the $L_X$ operators]
  \label{rem:pet-ch10-exercise-curvature-ricci-bracket-formula}
  \lean{PetersenLib.exercise10_5_26}
  \uses{thm:pet-ch10-curvature-bracket-formula}
  Let $M$ be a symmetric space and $X,Y,Z\in t_p$. Show
  $R(X,Y)Z=[L_X,L_Y]Z$ and $\operatorname{Ric}(X,Y)=-\operatorname{tr}([L_X,L_Y])$, where
  $L_X:=\operatorname{ad}_X$.
  \source{petersen:page-0407}
\end{remark}

\begin{remark}[Exercise 10.5.27 — parallel extension of a $p$-form via holonomy invariance]
  \label{rem:pet-ch10-exercise-parallel-form-holonomy-invariant}
  \lean{PetersenLib.exercise10_5_27}
  \uses{rem:pet-ch10-holonomy-properties}
  Consider a Riemannian manifold $M$ and a $p$-form $\omega$ on $T_pM$. Show
  that $\omega$ extends to a parallel form on $M$ if and only if $\omega$ is
  invariant under $\mathrm{Hol}_p(M)$.
  \source{petersen:page-0407}
\end{remark}
